% !TEX program = xelatex
\IfFileExists{Book\space NL/template/kaobook.cls}{
  \IfFileExists{main.tex}{\documentclass[main.tex]{subfiles}}{
    \IfFileExists{../../main.tex}{\documentclass[../../main.tex]{subfiles}}{
      \IfFileExists{../main.tex}{\documentclass[../main.tex]{subfiles}}{
        \documentclass[Book\space NL/main.tex]{subfiles}
      }
    }
  }
}{
  \IfFileExists{../../template/kaobook.cls}{\documentclass[../../main.tex]{subfiles}}{\documentclass{book}}
}
\IfFileExists{Book\space NL/template/kaobook.cls}{}{\IfFileExists{../../template/kaobook.cls}{}{\InputIfFileExists{review-notes-standalone.tex}{}{\InputIfFileExists{Book\space NL/appendices/appendix 4-Review Notes/review-notes-standalone.tex}{}{}}}}
\begin{document}
\ifSubfilesClassLoaded{
  \ifcsname appendixchapterparent\endcsname
  \else
    \standaloneappendixchaptersetup{D}{5}
  \fi
}{}
\IfFileExists{\booknlroot appendices/appendix 4-Review Notes/images/d5-figure-anatomy-pectoral-axilla-brachial-plexus-injuries.png}{
  \edef\reviewnotesimagepath{\booknlroot appendices/appendix 4-Review Notes/images/}
}{
  \edef\reviewnotesimagepath{\booknlroot appendices/appendix 4-Review Notes/images/}
}
\chapter{Integumentary and Musculoskeletal Systems}
\labch{appendix-d5-integumentary-musculoskeletal-review}
\margintoc
\noindent\textit{Transcribed from the supplied integumentary and musculoskeletal study notes. The content is grouped by topic rather than lecture number; spelling and grammar are lightly copyedited, while source claims are retained and accuracy concerns are flagged in bold where identified.}\par\medskip
% CONTENT AUDIT: The MD212 source lists L7, L17, L18, L22, and L27 as headings only, with no body notes; verify these remain intentionally unexpanded.

\section{Integument and Skin Disorders}

\subsection{Integumentary System}
\begin{itemize}
\item Skin characteristics: largest and heaviest organ (15\% body weight), 1-2 mm thick (most)
\item Skin functions:
\begin{itemize}
\item Protection from pathogens, chemicals, UV
\item Prevent loss of body fluids (by glycolipids at epidermis)
\item Regulation of temperature by sweat and blood vessel dilation/constriction
\item Sensory function (touch, pain, temperature, vibration and pressure)
\item Passive immune barrier (w/ langerhans immune cell = antigen presenting dendritic cell)
% CONTENT AUDIT: Retained the Word-source wording here; the precursor and UV conversion sequence should be checked against a reliable reference.
\item Vitamin D3 production from precursors (steroids (cholcalciferol) + UV light) \(\rightarrow\) used for calcium absorption from intestinal mucosa. \textbf{[Accuracy check: The source appears to misidentify the precursor and sequence; verify the role of 7-dehydrocholesterol, UVB, and cholecalciferol.] }
\end{itemize}
\item Skin is composed of many layers: (from top to bottom)
\begin{itemize}
\item Epidermis - Stratified squamous epithelium derived from the ectoderm \textbf{w/o vasculature}
\begin{itemize}
\item \emph{Stratum corneum }- top layer, barrier function. Absent nucleus, cytoplasm filled w/ keratin (by hydrolytic enzyme induced). Desmosome damaged \(\rightarrow\) corneocyte shedding (desquamation). Constantly replaced. Up to 30 layers thick, 20-30 days for full turnover (replacement).
\item \emph{Stratum lucidum }\textbf{- thick skin only - }translucent, dense. W/o nuclei and densely packed w/ eleidin (keratin intermediate). “Transition” between S. granulosum, S. corneum
\item \emph{Stratum granulosum} - 3-5 layers of flat keratinocytes. Contains:
\begin{itemize}
\item keratohyalin granules w/ cysteine rich content in lipid filled vesicles
\item lamellar granules w/ lipid to seal skin (make impermeable to H2O). Released into intercellular space
\item Nonmembranous flaggrin secretory granules. Flaggrin cross links w/ tonofilamen
\item NB: Keratin tonofilament+keratohyalin granules \(\rightarrow\) keratinization
\end{itemize}
\item \emph{Stratum spinosum }- 4-6 rows of keratinocytes appearing irregular+polyhedral+round nuclei w/ surface spikes = desmosome anchoring to intermediate/tonofilaments \(\rightarrow\) ‘prickle cell’ characteristic. S. spinosum synthesizes cytoplasmic intermediate filaments (cytokeratins) that aggregate to form tonofibrils. Tonofibrils bind to desmosomes and maintain cell cohesion/prevent abrasion. S. spinosum contains langerhans cells (modified macrophages)
\item \emph{Stratum basale} - basement membrane consisted of cuboidal/columnar stem cells (mitotically active) – precursor to keratinocytes. Produce and contain intermediate keratin filaments. Hemidesmosomes on basal surface, and desmosomes on the lateral surface
\end{itemize}
\item Dermis - dense fibrous connective tissue w/ blood vessels and nerves derived from the mesoderm (papillary and reticular layers)
\item Hypodermis - superficial fascia fascia w/ subcutaneous connective tissue and adipocytes
\end{itemize}
\item Attachment: Cell to extracellular matrix by hemidesmosome, and epidermis to dermis by collagen type IV
\item Skin is either thin or thick
\begin{itemize}
\item thick skin contains the stratum lucidum layer, which is not present in thin skin
\item Thick skin > 5mm thick, lacks sebaceous glands, hair. \textbf{[Accuracy check: Confirm the “>5 mm” threshold; thick skin is usually defined by epidermal thickness, while total skin thickness varies by site.] }
\item Covers palms and soles
\item Thick skin has more and taller dermal papilla
\end{itemize}
\item Cell types in epidermis:
\begin{itemize}
\item keratinocyte (majority) - produce keratin, originating in s. basale (produce keratin intermediate filaments, tonofilament – fibrous protein for protection). Tight intercellular bridge w/ desmosomes (Abnormal \(\rightarrow\) pemphigus). Continually more keratinized moving up to surface
\item Melanocyte - skin pigment granule production (sent to neighbor keratinocytes). Originates from neural crest derived cells. Round cell w/ dendritic processes. In the basal layer only. Protects against UV damage, free radical scavenging. Colour of skin =/= amount of melanocytes; colour is det/ by melanin granule amounts, speed of degradation for capsules. Appear as nucleus+clear cytoplasm in IHC
\item Langerhans cell - antigen presenting dendritic macrophage cell (present to CD8+ T cell). Responsible for immune hypersensitivity (contact allergic dermatitis). Indented nucleus and found in all layers (esp. S. spinosum). No desmosome, no tonofilament attachment to keratinocytes. Appear pale in H\&E
\item Merkel cells - Tactile sensory mechanoreceptors derived from neural crest. Associated w/ sensory nerve endings (merkel’s corpuscle). Found in S. basale and linked to keratinocytes by desmosomes. Numerous in fingertip, palmar and plantar (soles) skin.
\end{itemize}
\item Melanin pigment formation: Tyrosine \(\rightarrow\) DOPA \(\rightarrow\) Dopaquinone \(\rightarrow\) Phenomelanin, eumelain. Block tyrosinase/no reuptake of Tyr = no pigment = albinism
\item Skin colour factors: melanin accumulation in S. basale, S. spinosum (most significant) \(\rightarrow\) reddish brown to brown black, hemoglobin \(\rightarrow\) reddish to pinkish hue, carotene – yellow pigment from egg yolk, yellow/orange vegetables (accumulates in S. corneum, subcutaneous fat)
\item Dermis - Connective tissue layer of the skin.
\begin{itemize}
\item W/ elastic fibers to support epidermis, bind hypodermis
\item Assist in thermoregulation: sweat glands, blood vessels, sebaceous glands, nerve endings
\item Mainly collagen w. Elastic+reticular fibers+fibroblasts
\item Maintain elasticity, strength of skin (by collagen fibers, elastic, reticular fibers
\item Contains sensory receptors, nerve endings
\item Dermal - epidermal junction:
\begin{itemize}
\item Epidermal ridges (rete ridges) project from epidermis into papillary dermis (epidermal downgrowth). Impression of the ridges create fingerprint patterns – prevent shearing
\item Dermal ridges (dermal papillae) extends from dermis to epidermis (upwards)
\item Damage to junction = bullous pemphigoid (autoimmune blistering disorder caused by dermis-epidermis junction damage)
\end{itemize}
\item 2 main layers: papillary (20\%) and reticular (80\%)
\begin{itemize}
\item Papillary layer - less elastic/collagen fibers, more cells. Heavily vascularized (thermoregulation). W/ dermal, capillary loops and meisner’s corpsules. Dermal papillae extend into epidermis, forming fingerprints
\item Reticular layer - Dense irregular CT found in the deep dermis. Elastic w/ more collagen bundles and reticular fibers. Contain Glycoaminoglycans (hyaluronic acid, dermatan sulphate) for skin hydration
\end{itemize}
\end{itemize}
\item Hypodermis
\begin{itemize}
\item Superficial subcutaneous tissue above underlying tissues/organs. Consisted of loose connective tissue and adipose cells
\item Functionally stabilizes skin and is loosely attached to dermis and muscle. Contains sweat glands, hair follicles
\item Subcutaneous fat acts as an energy reservoir and thermal insulation. Is 8\% thicker in women
\item Highly vascular = site for drug introduction and quick absorption
\end{itemize}
\item Skin aging effects
\begin{itemize}
\item Epidermal thinning (w/ same no. of layers) \(\rightarrow\) skin appears more transparent
\item Melanocyte count decreases, hypertrophy \(\rightarrow\) age spots at sun exposed areas
\item Collagen and elastin changes reduce skin elasticity/strength \(\rightarrow\) wrinkles, fine lines, sagging skin
\item Sebaceous glands produce less oil \(\rightarrow\) skin dryness, itching
\item Thin subcutaneous fat layer \(\rightarrow\) reduced insulation, increased skin injury risk, less able to maintain temperature
\item Less sweating \(\rightarrow\) heat stroke risk
\end{itemize}
\item Clinical skin disorders (related to epidermis)
\begin{itemize}
\item \textbf{Psoriasis - Immune mediated inflammatory skin disorder w/ excessive melanocyte proliferation in S. basale and S. spinosum. Inflammatory cells, helper T cells, dendritic cells, neutrophils present. Elongation of epidermal papillae w/ angiogenesis. Gross presentation as sharply demarcated psoriatic plaque (white, scaly), pitting nails, oncholysis (nail-nailbed separation), salmon-patch, subungual hyperkeratosis (nail plate thickening)}
\item Pemphigus - Autoimmune disruption of desmosomes leading to weak intercellular connection
\item Squamous cell carcinoma \(\rightarrow\) second most common form of skin cancer, slow growing reddish or ulcerated legion originating from epidermal keratinocytes, w/ findings of enlarged cells and hyperchromatic nuclei. Found in sun exposed areas
\item Melanoma \(\rightarrow\) originates from melanocytes of S. basale. Most serious form
\item Basal cell carcinoma \(\rightarrow\) most common malignant cancer of skin. From the basal layer of epidermis, found in sun exposed areas. Appears as white nodules or waxy bumps on face/neck
\end{itemize}
\item Injuries to the skin
\begin{itemize}
\item Burns \(\rightarrow\) Demarcated by the number of layers affected - 1=epidermis, 2=dermis, 3=all, 4=including muscle
\item Cutaneous wounds (healing) - Bleeding into wound \(\rightarrow\) blood clotting, leukocyte cleaning \(\rightarrow\) vasculation, granulation tissue forms \(\rightarrow\) epithelial regeneration, connective tissue fibrosis
\end{itemize}
\item Epidermal derivatives (pilosebaceous apparatus):
\begin{itemize}
\item Hair:
\begin{itemize}
\item Hair follicle - tubular invagination w/ hair growth bulb \& dermal papilla
\item Hair shaft - cylindrical visible projection of hair (colour depends on amount, distribution of melanin)
\item Structure of hair follicle (outermost to innermost):
\begin{itemize}
\item First coat \(\rightarrow\) epidermal sheath (inner epithelial root sheath \(\rightarrow\) outer epithelial root sheath \(\rightarrow\) glassy membrane)
\item Second coat \(\rightarrow\) Connective tissue sheath - Surrounds outermost of hair. Attachment for arrector pili
\item Innermost edge surrounds the hair
\item NB: contains melanocytes for hair colour. Deep expanded part of hair follicle = hair bulb (made of pluripotent matrix cells). Differentiation, keratinization of hair bulb cells \(\rightarrow\) hair growth
\end{itemize}
\item Structure of hair (outermost below sheath to innermost):
\begin{itemize}
\item Cuticle - outermost layer of organized overlapping keratinous scales w/o nuclei
\item Cortex - Largest layer deeper than cuticle. Several layers of long, flattened, \textbf{spindle shaped cornified (hardened) cells}. Hard keratin fibrils run parallel to the long axis of hair. Pigment granules exist in this layer
\item Medulla - center of hair. Large vacuolated cells w/ shrunken and cornified cuboidal structure. Absent in fine short hair/languno hair. Soft keratin
\end{itemize}
\item Hair root???
\item Function: thermoregulation, heat conservation, sensation, body image
\item NB: skin of fetus is covered by lanugo (primary hair) later replaced by vellus (secondary hair)
\end{itemize}
\item Arrector pili muscle \(\rightarrow\) bundle of smooth muscle cells connecting to the papillary dermis. Function = increased protection against heat loss, goosebumps
\item Sebaceous glands (holochrine secretion) \(\rightarrow\) produce sebum, keeps hair and skin moist. Flask shaped w/ opening to the hair follicle. Consists of flattened epithelial cells on basement membrane w/ sebocytes (centrally placed round cells) filled w/ fat droplets. Stratified squamous duct. Not found in palm/sole, and controlled under adrenal androgen \(\rightarrow\) testosterone (male), ovarian hormone (female)
\end{itemize}
\item Epidermal derivatives:
\begin{itemize}
\item Sweat/cutaneous glands
\begin{itemize}
\item Apochrine sweat gland
\item \textbf{Merocrine sweat gland (aka eccrine)} \(\rightarrow\) produce sweat (water, sodium, ammonia, urea, uric acid). The secretory portion (acinus) is simple columnar/cuboidal epithelium w/: myoepithelial cell at basement membrane, clear cells (secrete water), dark cells (contain glycoprotein granules). Functions: thermoregulation, response to stress (cholinergic sympathetic nerve, Ach neurotransmitter). Found in both skin types (except lips, genitalia). Small lumen, not very eosinophilic
\item Merochrine duct \(\rightarrow\) simple coiled tubular glands w/ opening at epidermal surface (not associated w/ hair follicle). stratified cuboidal epithelium. Small cells and darker stains. Often intermixed w/ secretory units
\item Apocrine sweat glands \(\rightarrow\) a simple coiled tubular gland w/ secretory and myoepithelial cells. Found in axilla, areolar, nipple, circumanal regions (eg. genitalia). Thick and milky secretion (high fatty acid, protein). Can lead to rancid odor from bacteria in poor hygiene. Responds to sex hormones, emotional, sensory stimuli (not heat) by adrenergic nerves. Large lumen, eosinophilic
\item Apocrine duct \(\rightarrow\) only found connected to the hair follicle. Not intermixed w/ secretory units
\end{itemize}
\item Gland types:
\begin{itemize}
\item Merochrine - secreted by secretory vesicles. Found in eccrine gland (merochrine sweat gland)
\item Apocrine gland - pinching off parts of the cell to secrete
\item Holocrine - cell dying to secrete products by shedding off into duct
\end{itemize}
\item Nails
\begin{itemize}
\item Nail plate - keratinized part of S. corneum. Covers dorsal surface of terminal phalanges. Composed of nail body and nail root resting on nail bed (hard keratin, S. basale, S, spinosum)
\item Nail matrix - S. basale thickening near the nail root. Keratinocytes, melanocytes. Produces nail plate
\item Lunular - distal margin of generative epithelium of nail root
\item Lateral nail fold - skin of the lateral sides of the nail which is folded in
\item Systemic illness/malnutrition/chemotherapy \(\rightarrow\) onychomycosis (nail fungal infection)
\end{itemize}
\item Sensory nerve endings
\begin{itemize}
\item Non encapsulated - free nerve endings - pain, itching, temperature, touch (for around hair). No complex sensory structure. Found frequently in dermis, epidermis in hairy skin. Penetrate dermis and end in stratum granulosum, surrounding hair follicles. Penicillate fibers at the end of the nerve receive various inputs from epidermis/papillary dermis (except cold, which is sensed in papillary nerve ending)
\item Non encapsulated - Merkel’s corpuscle/disk - light touch
\item Encapsulated - Meissner’s corpuscle / tactile corpuscles - light touch sensors. Tapered cylinders arranged perpendicular to papillae and located inside of the dermis papillae. Each capsule = modified schwann cells. Most concentrated in thick hairless skin (esp. Finger pads). Sensitive to touch and point distinguishing.
\item \textbf{Encapsulated - Pacinian/lamellar corpuscle - mechanoreceptor responsible for sensing strong pressure / vibration. Large ovoid structure w/ concentric cellular lamellae. Found in deep dermis/hypodermis (esp. fingertips)}
\item Encapsulated - Ruffini corpuscle - continuous pressure
\end{itemize}
\end{itemize}
\item Clinical correlation (glands):
\begin{itemize}
\item Acne - increased sebum at puberty \(\rightarrow\) clogging follicles, irritation, bacterial colonization
\item Body odor - surface bacteria metabolize organic contents of sweat
\item Sweat - eccrine production may indicate disease - cystic fibrosis causes high NaCl in sweat (usually hypotonic), uremia in kidney failure causes urea in sweat
\end{itemize}
\item Circulation of the skin:
\begin{itemize}
\item Superficial plexus - between papillary and reticular dermis. Contains papillary loops
\item Cutaneous plexus - Deep, supplying hypodermis, hair follicles
\end{itemize}
\end{itemize}

\subsection{Skin Normal Flora}
\begin{itemize}
\item Skin structure (physical, chemical, physiological) and influence on microbiota community
\begin{itemize}
\item Physical factors - pH, temperature
\begin{itemize}
\item Skin pH typically below 5 - <5 is good, helps keep resident bacterial flora attached to skin (>8 would result in dispersal)
\item Skin temperature varies w/ environment (<37, about 30C)
\end{itemize}
\item Chemical factors - water, lipids, ions
\begin{itemize}
\item Salinity - sweat contains several ions (mainly Na and Cl). Microbes may be halotolerant/halophilic
\end{itemize}
\item Physiology - immunity
\begin{itemize}
\item Sweat has lysosome, antimicrobial peptide
\item Fatty acid in sebum is fungistatic
\item Dead tissue barrier
\item Antibodies
\end{itemize}
\end{itemize}
\item Predominant normal skin flora
\begin{itemize}
\item Microbiota - the assemblage of living organisms present in a defined environment
\item Normal flora = stable community of microbes permanently inhabiting various body sites in a healthy person (varies by person and body site). Typically symbiotic (mutualism, commensalism)
\item Includes microorganisms (bacteria, fungi, parasites, viruses)
\item Found in all areas of body exposed to environment
\item Are non-pathogenic at usual sites. Can become opportunistic pathogens at atypical sites. Some give host benefits eg. inhibition of pathogens, modulation of immune response (adaptive, innate immunity control)
\item Pathogen inhibition/blocking pathogenesis pathways:
\begin{itemize}
\item Byproduct release (eg. enzymes, bacteriocin, acid)
\item Competitive exclusion (nutrient, attachment site and area)
\item Enhance host innate immunity (AMP production, decreased inflammation, stronger epidermal barrier)
\item Enhance host adaptive immunity (local cytokine production tuning, epigenetically prime APCs to educate active immunity, regulatory T cell influence)
\end{itemize}
\item Can be either resident or transient flora (does it re-establish after disruption?)
\item Normal flora is not disrupted because: location of colonization (dead skin cells), low density and virulence, immunomodulatory molecule production, immune tolerance, co-evolution, other factors
\item Can include bacteria, viruses, fungi, parasites
\item Some areas are sterile sites (no microorganism growth)
\item Contributors to normal flora diversity:
\begin{itemize}
\item Transmission of non-resident microbes
\item Demographics
\item Host genetics
\item Human behavioural factors
\item Environmental characteristics
\end{itemize}
\end{itemize}
\item Bacteriome, microbiome, virome of skin \& effect of alteration
\begin{itemize}
\item Predominant resident microbes
\begin{itemize}
\item Nonhemolytic aerobic/anaerobic \emph{staphylococci} (eg. \emph{Staphylococcus epidermidis}, other coagulase -ve \emph{staphylococci}). Occasionally \emph{Staphylococcus aureus }and \emph{Peptidostreptococcus }speces
\item Alpha hemolytic streptococci (\emph{viridans streptococci})
\item \emph{Enterococcus }species (enterococci)
\item Aerobic/anaerobic \emph{diphtheroid} bacilli (eg. \emph{Corynebacterium}, \emph{Propionibacterium})
\item Gram positive aerobic spore forming \emph{bacilli} (air, water, soil)
\item Gram negative \emph{coliform }bacilli, \emph{Acinetobacter}
\item \emph{Canida sp.}, \emph{Malasseszia sp.}, \emph{Trichosporon sp.}
\end{itemize}
\item Viruses on skin
\begin{itemize}
\item Most = transient residents (usu. Due to contamination, transmission)
\item Some resident viruses (Human \emph{papillomavirus}, Merkel cell polyomavirus)
\item Primary immunodeficiency can find: \emph{Adenoviridae}, \emph{Anelloviridae}, \emph{Circoviridae}, \emph{Herpesviridae}, \emph{Papillomaviridae} and \emph{polyomaviridae}
\end{itemize}
\item Mites = mormal microbiota (eg. \emph{Demodex folliculorum})
\item Skin microbiome = the sum of microbes and their genomic elements in an environment (bacteriome, mycobiome, virome)
\begin{itemize}
\item Determination of microbiome - through 16s rRna (prokaryote) and 18s rRNA + internal transcribed spacer for eukaryotes
\end{itemize}
\item Alpha diversity = diversity between individuals in an environment
\item Beta diversity = diversity between samples in 2 or more systems
\item Actions altering skin microbiota: washing, bathing, usage of soap, food (gut-skin axis and skin excretion)
\end{itemize}
\item \textbf{Eubiosis and dysbiosis}
\begin{itemize}
\item \textbf{Eubiosis \(\rightarrow\) balance of normal microbiota community}
\item \textbf{Dysbiosis \(\rightarrow\) imbalance, disturbed eubiosis, can lead to many diseases (eg. acne, rosacea, atopic dermatitis, psoriasis vulgaris)}
\item \textbf{Treatment of dysbiosis: drugs, skin probiotics, microbiota transplantation}
\end{itemize}
\item Eye, ear normal flora
\begin{itemize}
\item Eye normal flora:
\begin{itemize}
\item Coagulase negative \emph{Staphylococcus}
\item \emph{S. aureus}
\item \emph{Streptococcus }species
\item \emph{Micrococcis}
\item Gram positive \emph{bacilli}
\item Eye has factors affecting bacteria: lysosome, IgA, temperature, blinking
\end{itemize}
\item Ear normal flora:
\begin{itemize}
\item \emph{Malasseszzia}
\item \emph{Staphylococcus}
\item \emph{Corynebacterium}
\item \emph{Bacillus}
\item \emph{Alloiococcus otitidis}
\item Composition of earwax: Long chain fatty acids (sat/unsat), alcohols, squalene, cholesterol, lysosome, antimicrobial peptides
\end{itemize}
\end{itemize}
\item Disadvantages of normal flora:
\begin{itemize}
\item Human body odor (from volatile compounds generated by bacteria)
\item Contamination of food, wound infection, opportunistic infections
\item Interfere w/ diagnosis
\item Nosocomial (healthcare associated infection), iatrogenic infection (physician acquired infection)
\end{itemize}
\end{itemize}
\begin{itemize}
\item NB: ไม่ออก disease based, อาจออก 2 types of staph (pathogenic variant), bacteria localized to skin
\end{itemize}

\subsection{Terminology of Clinical Dermatology}
*อาจารย์บอกว่าไม่ออกสอบ!

\begin{itemize}
\item Primary lesion = visible, unmodified
\item Secondary lesion = lesion that has been modified (eg. scratched)
\item Cardinal features of lesions: type, shape, arrangement, distribution
\item Skin biology
\item The epidermis has an average thickness of about 50 microns, and a turnover rate of 4 weeks
\item Inflammatory and immunologic disorders
\item Hair disorders
\item Vascular disorders
\item Traumatic and mechanical disorders
\item Cysts and skin tumors
\item Miscellaneous conditions
\end{itemize}

\subsection{Pathology of Skin}
\begin{itemize}
\item \textbf{Lichen planus -} presenting w/ buccal mucosa, wrist lesions.
\begin{itemize}
\item 5Ps:
\begin{itemize}
\item Pruritic
\item Purple
\item Polygonal
\item Planar
\item Papules and Plaques
\end{itemize}
\item Shiny surface and fine lines (Wickham striae)
\item Symmetrical distribution
\begin{itemize}
\item Wrist
\item Elbow
\item Extremities
\item Genitalia
\item Reticulated lace-like whitish patches
\end{itemize}
\item Pathogenesis \(\rightarrow\) CD8+ T cell–mediated cytotoxic immune response against Ag in the basal cell layer and the DEJ W/ unknown mechanisms
\item Treatment:
\begin{itemize}
\item May resolve spontaneously
\item Topical steroids
\item Immunosuppressive medication (in recalcitrant cases)
\end{itemize}
\end{itemize}
\end{itemize}
\begin{itemize}
\item Hair loss after surgery \(\rightarrow\) telogen effluvium (premature transition of hair follicles from anagen to telogen)
\item Bullous pemphigoid \(\rightarrow\) age>60 , tense blisters on normal/erythematous skin, marked pruritus. Commonly involving lower abdomen, thighs, flexors (rarely mucosa \(\rightarrow\) subdermal bullae. Eosinophils in Tzanck smear
\end{itemize}
\begin{center}
\captionof{table}{Comparison of pemphigus vulgaris and bullous pemphigoid.}\labtab{appendix-table-d5-pemphigus-bullous-pemphigoid}
\small
\renewcommand{\arraystretch}{1.15}
\begin{tabularx}{\textwidth}{p{0.22\textwidth}XX}
\toprule
\textbf{Feature} & \textbf{Pemphigus vulgaris} & \textbf{Bullous pemphigoid} \\
\midrule
Age & Young/middle age & Elderly \\

Severity & Fatal & Chronic and mild \\

Antibodies & Anti-desmoglein 1, 3 & -- \\

Skin & Flaccid blister & Tense blister \\

Mucosa & 90\% & Rare \\

Histopathology & Intraepidermal blister & Subepidermal blister \\

Tzanck smear & Acantholytic cell & Eosinophils \\

Treatment & High-dose immunosuppressive therapy & Topical or low-dose immunosuppressive therapy \\
\bottomrule
\end{tabularx}
\end{center}

\subsection{Dermatopharmacology}
\begin{itemize}
\item \textbf{Most penetration site \(\rightarrow\) mucus membrane}
\item \textbf{Fingertip unit = 0.5g (2 handprints). Face+neck = 2.5 FTU (1.25 g/day)}
\item \textbf{Mupirocin ointment (indicate in MRSA only, TH)}
\begin{itemize}
\item Use: superficial bacterial infection, intranasal for susceptible
\item individuals
\item Mechanism: inhibits protein synthesis by reversibly binding to bacterial isoleucyl transfer RNA synthetase
\item Active against staphylococcus groups A, B, C, G, streptococcus and some gram negative bacteria
\end{itemize}
\end{itemize}

\subsection{Wholistic Approach on Inflammatory and immunologic-Dermatology}
\begin{itemize}
\item \textbf{Psoriasis \(\rightarrow\) avoid alcohol consumption}
\item IgE receptors for urticaria w/ histamine involvement
\end{itemize}

\section{Regional Anatomy and Development}

\subsection{Anatomy of Pectoral, Axilla and Scapula}
\begin{itemize}
\item \textbf{Brachial plexus injuries}
\end{itemize}
\begin{center}\includegraphics[width=0.82\textwidth,height=0.68\textheight,keepaspectratio]{\reviewnotesimagepath d5-figure-anatomy-pectoral-axilla-brachial-plexus-injuries.png}\par\captionof{figure}{The illustration relates to Brachial plexus injuries, within the topic of Anatomy of Pectoral, Axilla and Scapula.}\labfig{appendix-figure-d5-anatomy-pectoral-axilla-brachial-plexus-injuries}\end{center}
\begin{itemize}
\item \textbf{Scapular winging - long thoracic nerve injury (C5 root if diaphragm involvement)}
\item \textbf{Shoulder joint abduction \(\rightarrow\) supraspinatus, deltoid}
\item \textbf{Intercostobrachial nerve supplies sensation to upper skin of axilla, medial upper arm (from T2 spine) \(\rightarrow\) common mastectomy injury}
\item \textbf{Dissect pectoralis major \(\rightarrow\) pectoral lymph node}
\end{itemize}

\subsection{Anatomy of Upper Limbs}
\begin{itemize}
\item Tennis elbow (cannot extend wrist) \(\rightarrow\) lateral epicondyle of humerus may be fractured
\item Lateral arm rotation weakness \(\rightarrow\) surgical neck of humerus fracture (axillary nerve injury)
\item Carpal tunnel syndrome \(\rightarrow\) median nerve compression
\end{itemize}

\subsection{Histology and Cellular Events of Muscle}
\begin{itemize}
\item Muscle tissue terminology:
\begin{itemize}
\item Muscle fiber - the functional and cellular unit of muscle tissue. Each fiber has a parallel array of cytoplasmic filament w/ myofilament
\item Myofilaments - The filaments which interact to induce muscle contractions. Two types, actin (thin filament) and myosin (thick filament)
\item Sarcolemma - cell membrane of muscle cell
\item Sarcoplasm - cytoplasm of muscle cell
\item Sarcoplastic reticulum - endoplasmic reticulum of muscle cell
\end{itemize}
\item Muscle tissue types:
\begin{itemize}
\item Skeletal muscle - produces purposeful movements of the skeleton (voluntary control). Peripherally located nucleus and multinucleated
\item Cardiac muscle - pumping action of the heart (involuntary control). Only type of muscle having intercalated disks. Centrally located nucleus. Has the most space between cells
\item Smooth muscle - motile force for the peristalsis of the gut and emptying of the urinary bladder. Found most in the hollow organs (rhythmic involuntary contractions)
\end{itemize}
\item Characteristics of muscle (skeletal, cardiac, smooth):
\begin{itemize}
\item Muscle properties:
\begin{itemize}
\item Irritability
\item Contractility
\item Extensibility
\item Elasticicity
\end{itemize}
\item Muscle functions:
\begin{itemize}
\item Movement
\item Heat generation
\item Posture
\item Joint stabilization
\end{itemize}
\item Skeletal muscle characteristics:
\begin{itemize}
\item sarcolemma (cell membrane) \(\rightarrow\) endomysium \(\rightarrow\) epimysium \(\rightarrow\) perimysium
\item Epimysium \(\rightarrow\) dense irregular collagenous CT, surrounds entire muscle
\item Perimysium \(\rightarrow\) CT surrounding each of the fascicles (muscle units)
\item Endomysium \(\rightarrow\) surrounds each muscle cell/muscle fiber. Contains satellite cells and capillaries/nerve fibers
\item At the ends of the muscle, the three layers come together and form CT attachment to the bone matrix, either bundles (tendons) or sheet (aponeurosis)
\item Long muscle fibers (=skeletal muscle cell) and large size (originally from fusion of mesodermal cells, from myoblasts)
\item Multinucleated (hundreds of nuclei)
\item Contains satellite (stem) cells to repair damage
\item Organization of skeletal muscle: myofilaments (in cells) \(\rightarrow\) myofibril (sarcomere chain) \(\rightarrow\) muscle fiber/cell \(\rightarrow\) fascicle (bundle of muscle fibers) \(\rightarrow\) muscle (many bundles of fiber)
\item Structure under microscope: long, cylindrical, multinucleated cell w/ nuclei peripherally located. Surrounded by endomysium. Orderly arrangement of myofibrils showing light and dark cross striations w/ alternating light (l) and dark (A) bands
\item Transverse Z disks mark the end of each sarcomere (anchor thin/actin filament). The center of sarcomere has thick (myosin) filament, forming A bands. The band is formed by thin filaments. The H zone (cross striation) is formed by thick filament, and the M line is the middle of each sarcomere
\item Each sarcomere extends from Z disk to the next Z disk
\end{itemize}
\end{itemize}
\end{itemize}
\begin{center}\includegraphics[width=0.82\textwidth,height=0.68\textheight,keepaspectratio]{\reviewnotesimagepath d5-figure-histology-cellular-events-each-sarcomere-extends-disk.png}\par\captionof{figure}{The illustration relates to Each sarcomere extends from Z disk to the next Z disk, within the topic of Histology and Cellular Events of Muscle.}\labfig{appendix-figure-d5-histology-cellular-events-each-sarcomere-extends-disk}\end{center}
\begin{itemize}
\item Major components of skeletal muscle fibers
\begin{itemize}
\item Sarcolemma, sarcoplasma, sarcoplasmic reticulum, transverse tubule (where the cell membrane extends into the sarcoplasm, invagination), terminal cisternae (site for Ca2+ storage), myofibrils in sarcomeres, muscle triads (T tubules + 2 cisternae)
\end{itemize}
\item Each myofibril (part of muscle cell) consists of 3 protein types: contractile, regulatory, structural:
\begin{itemize}
\item Contractile proteins \(\rightarrow\) actin (thin filament), myosin (thick filament)
\item Myosin \(\rightarrow\) many myosin molecules per each actin filament. Golf club shape (thickening) and long tail. Actin binding site on the head. ATPase at the head releases energy for muscle contraction
\item Actin \(\rightarrow\) formed from 3 proteins globular (G) to bind to myosin, tropomyosin (long fibrous molecule, prevents acin-myosin binding) and troponin (reversible binding in presence of Ca2+ to move tropomyosin from active site)
\item Regulatory proteins \(\rightarrow\) tropomyosin, troponin
\item Structural proteins \(\rightarrow\) titin, nebulin, dystrophin
\end{itemize}
\end{itemize}
\begin{itemize}
\item Components of thin/thick filaments of sarcomere of skeletal muscle
\item Steps involved in 3 sliding mechanism of skeletal muscle contraction
\item Neuromuscular junction
\begin{itemize}
\item Action potential \(\rightarrow\) momentary change in cell electrical potential caused by ion influx/efflux occurring when a cell is stimulated. Transmits a nerve impulse
\item Neuromuscular junction \(\rightarrow\) Junction between motor neuron and axon terminal. Located at a highly excitable region of muscle membrane, and initiates action potential responses. The NMJ junction uses acetylcholine in vertebrates
\item The NMJ contains: presynaptic terminal, the sarcolemma, synaptic vesicles, nicotinic acetylcholine receptors and mitochondria (at the sarcoplasm side)
\item Process of signal transmission at the NMJ:
\begin{itemize}
\item Acetylcholine (ACh) released from axon terminal of motor neuron
\item ACh binds to receptors in the motor end plate. End plate potential is triggered, resulting in an action potential
\item Action potential propagation through sarcolemma and down T tubules
\item Ca2+ is released from the sarcoplasmic reticulum
\item Ca2+ binds to troponin causing conformation change in tropomyosin which exposes myosin binding sites
\item Crossbridging of actin begins
\end{itemize}
\end{itemize}
\item Cardiac muscle characteristics:
\begin{itemize}
\item Striated w/ same contractile machinery as skeletal muscle. However, is self excitatory (autonomically modulated)
\item Innervation is neuroendocrine (no motor end plates)
\item Features centrally placed nuclei (often w/ biconical clear area next to nucleus) and branched fibers joined by intercalated disks
\item Many mitochondria (up to 40\% cell volume)
\item Sarcoplasmic reticulum+T-tubules in dyads at Z lines
\item Shorter terminal cisternae of the sarcoplasmic reticulum
\item Tubules 2x larger in diameter than in skeletal muscle (to aid Ca2+ transport)
\end{itemize}
\item Smooth muscle characteristics:
\begin{itemize}
\item Found in most blood vessel walls, GI and urogenital organ walls, arrestor pili muscle
\item Each fiber is fusiform w/ centrally located nucleus
\item Cytoplasm has both actin and thick filament. No striation or overlapping
\item Fibers are connected by gap junctions (involuntary autorhythmicity)
\item Contractions occur in a neuroendocrine manner
\end{itemize}
\item See the histology slides of LS and TS for cardiac, smooth and skeletal muscle
\end{itemize}

\subsection{Histology and Cellular Events of Cartilage and Bone}
\subsubsection{Cartilage Tissue}
\begin{itemize}
\item Functional unit of joints: cartilage, ligaments, and synovium
\item Macroscopic and microscopic of cartilage:
\begin{itemize}
\item Connects between two bones \(\rightarrow\) firm and resilient CT (bear stress w/o permanent distortion). Avascular tissue. Supports soft tissues
\item Made of cells (fibroblasts and chondrocytes), extracellular matrix (made of type II collagen) and ground substance (Hyaluronic acid + glycosaminoglycans)
\item Can act as an embryonic precursor/model for long bones. Is all replaced w/ bone except in articular surfaces
\end{itemize}
\item Growth of cartilage
\begin{itemize}
\item Chondroblasts \(\rightarrow\) chondrocyte progenitor. Lines border between perichondrium and matrix
\item Chondroblasts - function in chondrogenesis - deposit type II collagen and extracellular matrix (hyaluronic acid, glycosaminoglycans – mainly chondroitin sulfate and keratin sulfate)
\item Chondrocyte \(\rightarrow\) mature cartilage cells reside in and fill the lacuna (a space in osteons). Produces type II collagen and ECM (hyaluronic acid, GAGs, mainly chondroitin sulfate and keratin sulfate)
\item Chondrogenesis (histogenesis of cartilage)
\begin{itemize}
\item Initially - matrix secretion
\item Continued w/ appositional growth/interstitial growth
\end{itemize}
\item \textbf{Appositional growth} \(\rightarrow\) growth from perichondrium (outer fibrous layer, inner cellular layer) \(\rightarrow\) progenitor cells in perichondrium produces new cartilage from outside
\begin{itemize}
\item Cartilage added just underneath perichondrium
\item Chondroblasts divide into immature chondrocytes, which secrete new matrix
\item Most common in cartilage growth
\item Perichondrium layer differentiation, type II collagen expression and proteoglycan aggregation controlled by Sox9
\end{itemize}
\item Interstitial growth \(\rightarrow\) growth from within cartilage \(\rightarrow\) chondrocytes increase the number of cells and amount of matrix
\begin{itemize}
\item Occurs in early stage of hyaline cartilage (joint cartilage) formation
\item Is the formation of new cartilage from within the cartilage mass
\item Present in articular cartilage (joints) due to lack of pericondrium
\item very active during endochondral ossification (when cartilage is being replaced by bone)
\end{itemize}
\end{itemize}
\end{itemize}
\begin{center}\includegraphics[width=0.82\textwidth,height=0.68\textheight,keepaspectratio]{\reviewnotesimagepath d5-figure-histology-cellular-events-very-active-endochondral-ossification.png}\par\captionof{figure}{The illustration relates to very active during endochondral ossification (when cartilage is being replaced by bone), within the topic of Histology and Cellular Events of Cartilage and Bone.}\labfig{appendix-figure-d5-histology-cellular-events-very-active-endochondral-ossification}\end{center}
\begin{center}\includegraphics[width=0.82\textwidth,height=0.68\textheight,keepaspectratio]{\reviewnotesimagepath d5-figure-histology-cellular-events-histology-cellular-events-cartilage.png}\par\captionof{figure}{The illustration relates to Histology and Cellular Events of Cartilage and Bone, within the topic of Histology and Cellular Events of Cartilage and Bone.}\labfig{appendix-figure-d5-histology-cellular-events-histology-cellular-events-cartilage}\end{center}
\begin{itemize}
\item Types of cartilage
\begin{itemize}
\item Hyaline cartilage (Type II collagen) \(\rightarrow\) articular surfaces, Epiphyseal plate, tracheal wall,cartilage of respiratory tract
\begin{itemize}
\item Contains type II collagen, glycosaminoglycans, proteoglycans
\item Has a perichondrium (except articular cartilage) which forms a capsule around the cartilage
\item High water content and mobility allow for weight bearing capacity
\item Normal matrix turnover \(\rightarrow\) cells respond to changes in matrix composition
\item Pressure loads signal for chondrocyte synthesis activity (mechanical, electrical, chemical signals)
\item As aging occurs, chondrocytes are less able to respond to matrix changes
\item Superficial zone (elongated, flattened chondrocytes) \(\rightarrow\) intermediate zone (round chondrocytes) \(\rightarrow\) deep zone (short columns of chondrocytes) \(\rightarrow\) calcified zone (borders bone, calcified matrix w/ chondrocytes)
\end{itemize}
\item Elastic cartilage (Type II collagen, elastic fibers) \(\rightarrow\) pinna of the ear, epiglottis, eustachian tube
\begin{itemize}
\item Hyaline cartilage w/ elastic fibers. Has both firm and elastic properties
\item Has a perichondrium
\item No calcification present
\end{itemize}
\item Fibrocartilage (Type II and Type I collagen) \(\rightarrow\) intervertebral disks, TMJ joint, pubic symphysis, insertion sites of tendons and ligaments
\begin{itemize}
\item Hyaline cartilage + dense regular connective tissue
\item High in collagen type I and II
\item No perichondrium
\item Resists compression, shearing, is a shock absorber
\end{itemize}
\end{itemize}
\item Cartilage repair:
\begin{itemize}
\item Limited repair \(\rightarrow\) chondrocytes cannot move, limited proliferation. Typically repaired from perichondrium
\end{itemize}
\item Hyaline calcification \(\rightarrow\) used to replace cartilage w/ bone. Calcifies over time as humans age leading to reduced joint mobility
\item \textbf{Things that represent fibro/elastic cartilage}
\item \textbf{Exceptions: what is not hyaline cartilage, where is it located?}
\item \textbf{Type II collagen?}
\end{itemize}
\subsubsection{Bone}
\begin{itemize}
\item Overview \(\rightarrow\) bone is made of bone cells, bone matrix, periosteum. Endosteum
\begin{itemize}
\item Bone supports the body, protects vital organs. cranial, thoracic cavities. Also contains bone marrow
\item Act as a series of levers which multiply force from muscle contractions and transform into body movements
\item Bone functions as a reservoir for: calcium, phosphate, other ions
\end{itemize}
\item Cells of bone
\begin{itemize}
\item Osteoblasts
\begin{itemize}
\item Synthesizes and secretes organic components of bone matrix
\item Consists of type I collagen, proteoglycan, glycoprotein
\item Synthesizes organic components of bone matrix and usually located at outermost surface of bone matrix
\item Differentiate from osteoprogenitor cells in periosteum (outer layer of bone) and endosteum (outer layer of cancellous bone)
\item cover the surfaces of existing bone matrix called “bone lining cells”
\item Secretes osteoid (contains collagen type I, proteoglycans, glycoproteins such as osteonectin (hardens to form bone) and also secretes matrix vesicles)
\end{itemize}
\end{itemize}
\end{itemize}
\begin{center}\includegraphics[width=0.82\textwidth,height=0.68\textheight,keepaspectratio]{\reviewnotesimagepath d5-figure-histology-cellular-events-secretes-osteoid-contains-collagen.png}\par\captionof{figure}{The illustration relates to Secretes osteoid (contains collagen type I, proteoglycans, glycoproteins such as osteonectin (hardens to form bone) and also secretes matrix vesicles), within the topic of Histology and Cellular Events of Cartilage and Bone.}\labfig{appendix-figure-d5-histology-cellular-events-secretes-osteoid-contains-collagen}\end{center}
\begin{itemize}
\item Can differentiate into osteocytes
\item Osteocytes
\begin{itemize}
\item Found in cavity (lacuna) between bone matrix layers w/ cytoplasmic process extending into small canaliculi for communication channels (in gap junctions). Also used for nutrient transport
\item Gradually surrounded by their secretions
\item Eventually enclosed into lacunae
\item Is a mature bone cell enclosed in bone matrix
\item Extends dendritic processes to other osteocytes to communicate
\item Maintain calcium homeostasis, act as biosensors for detecting mechanical stress on bone (direct bone remodelling)
\item Responsible for maintenance and turnover of bone matrix
\end{itemize}
\item Osteoclasts - multinucleated giant cells involved in the resorption and remodelling of bone
\begin{itemize}
\item Responsible for bone resorption
\item Multinucleated cells at sites of bone removal
\item Resorption bay/Howship’s lacuna = shallow bay of bone
\item Derived from hematopoietic cells, granulocyte/macrophage progenitor cells
\item When resorbing bone \(\rightarrow\) ruffled border (cytoplasmic projections contacting bone) w/ microvilli like structure of membrane
\item Clear zone (sealing zone) ring like structure of cytoplasm adjacent to ruffled border. Contains cells, ecm, adhesion molecule, etc. Integrin, vinculin and talin
\item Basolateral region responsible for exocytosis of digested material
\end{itemize}
\end{itemize}
\begin{center}\includegraphics[width=0.82\textwidth,height=0.68\textheight,keepaspectratio]{\reviewnotesimagepath d5-figure-histology-cellular-events-basolateral-region-responsible-exocytosis.png}\par\captionof{figure}{The illustration relates to Basolateral region responsible for exocytosis of digested material, within the topic of Histology and Cellular Events of Cartilage and Bone.}\labfig{appendix-figure-d5-histology-cellular-events-basolateral-region-responsible-exocytosis}\end{center}
\begin{itemize}
\item Types of bone
\begin{itemize}
\item Lamellar bone
\begin{itemize}
\item Majority of bone in the body
\item Main: multiple layers (lamilae) of calcified matrix
\item Osteon/haversian systems \(\rightarrow\) Made of complex and concentric lamellae surrounding a small lamellar canal (w/ blood vessels, nerves, loose connective tissue and endosteum)
\item Volkmann canal (traverse canal) - narrow transverse canal for central canal to communicate w/ marrow cavity
\item Outer border of each osteon = cement line (collagen rich)
\item Organization of lamellae - external circumferential lamellae \(\rightarrow\) interstitial lamellae \(\rightarrow\) internal circumferential lamellae (around marrow cavity (is a part of the osteon)
\item Osteon - Consisted of the area surrounded by interstitial lamellae w/ the haversian canal inside (osteocytes, cytoplasmic processes, haversian canal)
\end{itemize}
\item Woven bone
\begin{itemize}
\item Non osteon/haversian system
\item Random deposition of type I collagen - appears in embryonic development and in fracture repair
\item Temporary - replaced by lamellar bone (except near skullcap, tendon insertions). Higher \% of osteocytes than bone
\end{itemize}
\end{itemize}
\item Gross observations of bone:
\begin{itemize}
\item Compact (cortical) bone 80\%
\begin{itemize}
\item Haversian system or osteon
\item Concentric Lamella
\item Haversian canal
\item Osteocyte in lacuna
\item Canaliculi (cytoplasmic process)
\item Volkman’s canal (Transverse perforating canal)
\end{itemize}
\item Cancellous (trabecular or spongy) bone 20\%
\begin{itemize}
\item Trabecula or Spicule
\item Concentric lamella
\item Non Haversian system
\end{itemize}
\end{itemize}
\item Bone classification by shape:
\begin{itemize}
\item Long bones: epiphysis (expanding end of long bone, spongy bone covered by a thin amount of compact bone) \(\rightarrow\) diaphysis (main shaft of bone, compact bone w/ a central marrow cavity) \(\rightarrow\) epiphyseal plate (hyaline cartilage between epiphysis and diaphysis)
\item Short bones, flat bones, irregular bones
\end{itemize}
\item Bone marrow:
\begin{itemize}
\item White bone marrow - fatty bone marrow, replaces red bone marrow gradually as people age. Contains an abundance of adipose cells. Does not produce blood cells. Found in the diaphyseal portion of long bones. During anemia can turn back to red bone marrow and produce blood
\item Red bone marrow - in children, is the content of every medullary cavity and contains an abundance of red blood cells and hematopoietic stem cells. Limited to skull, vertebrae, sternum, ribs, vertebrae, sternum, ribs, part of the pelvic (hip) girdle and proximal head of humerus and femur
\end{itemize}
\item Osteogenesis
\begin{itemize}
\item Intramembranous ossification
\item Endochondral ossification
\end{itemize}
\item Bone remodelling, repair
\end{itemize}

\subsection{Anatomy of Gluteal and Thigh}
\begin{itemize}
\item \textbf{Suspension of pelvis during one leg \(\rightarrow\) superior gluteal nerve}
\item \textbf{Extracapsular femur fracture \(\rightarrow\) medial circumflex femoral artery most likely injury}
\end{itemize}

\subsection{Development of Musculoskeletal System}
\begin{itemize}
\item \textbf{Most of the mesenchyme in the head is derived from the neural crest (musculature of the face and neck), facial bone}
\item \textbf{apical ectodermal ridge [AER] = induces finger development and order through chemical concentration gradient}
\item \textbf{Spina biffida - failure of the halves of the embryonic cartilaginous neural arch to fuse results in various types of spina bifida}
\end{itemize}
\begin{center}\includegraphics[width=0.82\textwidth,height=0.68\textheight,keepaspectratio]{\reviewnotesimagepath d5-figure-development-musculoskeletal-spina-biffida-failure-halves.png}\par\captionof{figure}{The illustration relates to Spina biffida - failure of the halves of the embryonic cartilaginous neural arch to fuse results in various types of spina bifida, within the topic of Development of Musculoskeletal System.}\labfig{appendix-figure-d5-development-musculoskeletal-spina-biffida-failure-halves}\end{center}
\begin{center}\includegraphics[width=0.82\textwidth,height=0.68\textheight,keepaspectratio]{\reviewnotesimagepath d5-figure-development-musculoskeletal-development-musculoskeletal-system.png}\par\captionof{figure}{The illustration relates to Development of Musculoskeletal System, within the topic of Development of Musculoskeletal System.}\labfig{appendix-figure-d5-development-musculoskeletal-development-musculoskeletal-system}\end{center}
\begin{center}\includegraphics[width=0.82\textwidth,height=0.68\textheight,keepaspectratio]{\reviewnotesimagepath d5-figure-development-musculoskeletal-development-musculoskeletal-system-02.png}\par\captionof{figure}{The illustration relates to Development of Musculoskeletal System, within the topic of Development of Musculoskeletal System.}\labfig{appendix-figure-d5-development-musculoskeletal-development-musculoskeletal-system-02}\end{center}
\begin{center}\includegraphics[width=0.82\textwidth,height=0.68\textheight,keepaspectratio]{\reviewnotesimagepath d5-figure-development-musculoskeletal-development-musculoskeletal-system-03.png}\par\captionof{figure}{The illustration relates to Development of Musculoskeletal System, within the topic of Development of Musculoskeletal System.}\labfig{appendix-figure-d5-development-musculoskeletal-development-musculoskeletal-system-03}\end{center}
\begin{center}\includegraphics[width=0.82\textwidth,height=0.68\textheight,keepaspectratio]{\reviewnotesimagepath d5-figure-development-musculoskeletal-development-musculoskeletal-system-04.png}\par\captionof{figure}{The illustration relates to Development of Musculoskeletal System, within the topic of Development of Musculoskeletal System.}\labfig{appendix-figure-d5-development-musculoskeletal-development-musculoskeletal-system-04}\end{center}
\begin{center}
\captionof{table}{Common structural types of limb malformations.}\labtab{appendix-table-d5-limb-malformations}
\small
\renewcommand{\arraystretch}{1.15}
\begin{tabularx}{\textwidth}{p{0.30\textwidth}X}
\toprule
\textbf{Term} & \textbf{Description} \\
\midrule
Amelia (ectromelia) & Absence of an entire limb \\

Acheiria, apodia & Absence of hands or feet \\

Phocomelia & Absence or shortening of proximal limb segments \\

Hemimelia & Absence of preaxial or postaxial parts of a limb \\

Meromelia & General term for absence of part of a limb \\

Ectrodactyly & Absence of any number of digits \\

Polydactyly & Excessive number of digits \\

Syndactyly & Presence of interdigital webbing \\

Brachydactyly & Shortened digits \\

Split hand or foot & Absence of central components of hand or foot \\
\bottomrule
\end{tabularx}
\end{center}
\begin{center}\includegraphics[width=0.82\textwidth,height=0.68\textheight,keepaspectratio]{\reviewnotesimagepath d5-figure-development-musculoskeletal-development-musculoskeletal-system-05.png}\par\captionof{figure}{Clinical examples of limb malformations.}\labfig{appendix-figure-d5-development-musculoskeletal-development-musculoskeletal-system-05}\end{center}
Pathology Labs

See Anki deck: https://ankiweb.net/shared/info/1658159405?cb=1760620150240

On the summative:

Osteosarcoma - Malignant matrix forming bone tumor (given picture of histology)

Veruca vulgaris (from HIV) - Identify koilocytes (given picture of histology). True false questions

???

SUM2


\section{Muscle, Joint, and Movement}

\subsection{Muscle Physiology}
\begin{itemize}
\item น่าจะออกเป็น mechanism of contraction for each muscle type
\end{itemize}
\subsubsection{Skeletal Muscle Physiology}
\begin{itemize}
\item \textbf{Motor unit recruitment size principle - slow and fast twitch fibers recruited in relation to exercise. As the effort becomes more intense, progressively more fast-twitch fibers are recruited}
\item \textbf{Types of muscle contraction:}
\begin{itemize}
\item \textbf{Isotonic: concentric contraction \(\rightarrow\) muscle shortens. Force > resistance}
\item \textbf{Isotonic: eccentric contraction \(\rightarrow\) muscle lengthens. Force < resistance}
\item \textbf{Isometric (static) contraction \(\rightarrow\) no movement \(\rightarrow\) force = resistance}
\end{itemize}
\end{itemize}
\begin{itemize}
\item \textbf{Types of muscle fiber:}
\begin{itemize}
\item Slow oxidative \(\rightarrow\) responsible for sustained movements, weight bearing antigravity
\item Fast oxidative \(\rightarrow\) sustained locomotion
\item Fast glycolytic \(\rightarrow\) responsible for fast ‘burst’ movements
\end{itemize}
\item \textbf{Length-tension relationship \(\rightarrow\) }filament overlap effects force generation of muscle
\end{itemize}
\begin{center}\includegraphics[width=0.82\textwidth,height=0.68\textheight,keepaspectratio]{\reviewnotesimagepath d5-figure-muscle-physiology-length-tension-relationship-filament.png}\par\captionof{figure}{The illustration relates to Length-tension relationship \(\rightarrow\) filament overlap effects force generation of muscle, within the topic of Muscle Physiology.}\labfig{appendix-figure-d5-muscle-physiology-length-tension-relationship-filament}\end{center}
\begin{itemize}
\item Muscle adaptations:
\begin{itemize}
\item \textbf{Muscle hypertrophy} \(\rightarrow\) Contractile elements enlarge (more sarcomeres in parallel) - occurs in weight training or overload
\item \textbf{Muscle hyperplasia} \(\rightarrow\) Increased number of fiber within a muscle (more sarcomeres in parallel) - Unlikely to occur in humans
\item Lengthening \(\rightarrow\) Occur during growth and eccentric-only weight training (muscle stretch training)
\end{itemize}
\item \textbf{Neuromuscular junction disorders}
\begin{itemize}
\item Botulinum toxins \(\rightarrow\) Prevents fusion of neurotransmitter vesicles w/ neuron cell membrane, leading to a lack of contraction signalling in muscle tissue
\item Myasthenia gravis \(\rightarrow\) Skeletal muscle receptor blocking by anti-AchR antibody leading to paralysis
\item Lambert‐Eaton myasthenic syndrome \(\rightarrow\) antibodies against presynaptic calcium channels at the motor nerve terminal of the NMJ decrease Ach release and results in weakness
\end{itemize}
\end{itemize}
\subsubsection{Smooth Muscle Physiology}
\begin{itemize}
\item Characteristics
\begin{itemize}
\item \textbf{Reverse stress-relaxation: initially increases tension when first subject to stress. Followed by relaxation of the muscle, reduced tension (eg. in urinary bladder)}
\item Located in wall of visceral organs, vessels, ciliary muscle, pupillary muscle, arrector pili
\item Under involuntary control
\item Present w/ caveolae instead of T tubules (small membrane invaginations)
\item Intermediate filaments link dense bodies into a cytoskeletal network
\item No striations in muscle tissue (contractile filaments not in alignment)
\item Contain my tropomyosin but no troponin (thick filament regulated)
\item Actin myosin in a 10:1 ratio
\item \textbf{Slow to contract because slow ATPases in smooth muscle lead to a decreased rate of crossbridge cycling}
\end{itemize}
\item Single unit vs multi unit muscle fibers:
\begin{itemize}
\item Single unit muscle fibers
\begin{itemize}
\item sparse innervation by both peripheral and sympathetic nervous system
\item Spontaneous and autorhythmic contractions (by pacemaker cells in GI tract only)
\item Regulated by nerve impulse
\item Active tension when stretched
\item Gap junction/nexi facilitate conduction of electrical activity between fibers
\item Found in GI tract, uterus, urinary bladder, small arterioles
\end{itemize}
\item Multi unit muscle fiber
\begin{itemize}
\item More concentrated innervation (than single unit smooth muscles)
\item Singularily innervated (by PNS or autonomic nervous system and varicose veins)
\item Lack spontaneous contraction, no stress response
\item Found in: ciliary muscle, iris muscles, vas deferens, GI sphincters, larger blood vessels
\end{itemize}
\end{itemize}
\item Control of smooth muscle activity
\begin{itemize}
\item Types of smooth muscle (based on activity pattern)
\begin{itemize}
\item Phasic smooth muscle - smooth muscle which acts rhythmically (eg. peristalsis in the GI / urogenital tracts)
\item Tonic smooth muscle - maintains a constant tone, constantly contracted to some extent (eg. vascular smooth muscle, bronchial, some sphincters)
\end{itemize}
\item \textbf{Smooth muscles can develop force rapidly but can slow down ATPase, resulting in maintenance of contractile force (usually 20-30\% of cross bridges). This is referred to as the latched state}
\item Smooth muscles depolarize w/ L type voltage gated Ca2+ channels, repolarize w/ K+ channels
\item Action potential of smooth muscle can be prolonged or brief
\item Some smooth muscles spontaneously generate slow waves (contraction-relax in a wavelike manner) or rhythmic contractions
\item Some smooth muscles contract based on hormone levels w/o changing action potential (latch state muscles)
\end{itemize}
\item Regulation of myoplasmic [Ca2+]
\begin{itemize}
\item Both extracellular and intracellular Ca2+ are able to initiate a contraction
\item Mechanisms to initiate Ca2+ release
\begin{itemize}
\item Voltage gated Ca2+ channels
\item Ligand gated Ca2+ channels
\item Mechanically gated / stretch activated Ca2+ channels \(\rightarrow\) eg. blood vessels on distortion of cell membrane, myogenic (from muscle tissue) contraction in the urinary bladder
\end{itemize}
\item Many mechanisms cause Ca2+ efflux from the sarcoplasmic reticulum, including IP3 gated Ca2+ channels, Gq activation (which opens gated channel or Ryanodine receptors)
\item Ca2+ enters cell through caveoli (in folding of membrane) channel
\item Ca2+ releases from sarcoplasmic reticulum during Ca2+ induced Ca2+ release or IP3 activation of sarcoplasmic Ca2+ channels
\item Ca2+ influx by carcolemma includes from leak Ca2+ channels, voltage gated, ligand gated and stretch activated non specific channels
\end{itemize}
\item Regulation of contraction
\begin{itemize}
\item Smooth muscle may have sensitivity to one or more of: spontaneous electrical activity, neurotransmitters, hormones, local changes in chemical composition and stretching
\item Intestinal smooth muscle membrane potential during contraction
\begin{itemize}
\item In the resting stage, slow waves occur
\item During stimulation (by stretching, acetylcholine or parasympathetics), spikes occur on the slow waves in addition to depolarization
\item When stimulated by the sympathetic nervous system or norepinephrine, the resting membrane potential decreases
\end{itemize}
\item Thick filament regulation entails the cross bridge cycle being controlled by phosphorylation of myosin light chain by myosin light chain kinase
\item Initiation of cross bridge cycling:
\begin{itemize}
\item Calcimodulin binds w/ Calcium to form Ca2+ - calcimodulin complex
\item Ca2+ - calcimodulin activates myosin light chain kinase
\item Myosin light chain kinase phosphorelates regulatory light chain of myosin to permit actin binding \(\leftarrow\) ‘thick filament regulated contraction’
\item Crossbridge cycle generates force of contraction
\item Decreased intracellular Ca2+ leads to myosin light chain dephosphorylation by myosin phosphatase
\item Relaxation or sustained contraction (because of latched bridge mechanism)
\end{itemize}
\item \textbf{Contraction or inhibition of contraction can be initiated by:}
\begin{itemize}
\item \textbf{Intrinsic action of pacemaker cells (in GI tissue)}
\item \textbf{Neurally released transmitters}
\item \textbf{Circulating / locally generated hormones and signalling molecules}
\end{itemize}
\item \textbf{Excitation-contraction coupling in smooth muscle}
\item Maintenance of latch state:
\begin{itemize}
\item Phosphorylated crossbridges can dephosphorylate while still attached to actin \(\rightarrow\) rigor like state (isometric)
\item Slow crossbridge cycling \(\rightarrow\) low ATP consumption per force exerted \(\rightarrow\) can maintain tension for long durations w/ low ATP consumption
\item Tonic contraction at low energy expenditure = latch state
\end{itemize}
\end{itemize}
\item Biophysical properties
\begin{itemize}
\item Length-tension relationship
\begin{itemize}
\item Smooth muscle is able to generate near maximum force at a range of lengths (broader range than skeletal muscle)
\item Overall length-tension relationship remains similar to skeletal muscle
\item The wider range of optimal lengths is hypothesized to be from increasing the number of units in series
\item Muscle strength-relaxation curve for smooth muscle tension shows plasticity when experiencing distension
\end{itemize}
\item Force-velocity relationship
\begin{itemize}
\item Low ATPase activity of myosin isoform in smooth muscle \(\rightarrow\) maximum speed of smooth muscle is slower than skeletal muscle
\end{itemize}
\end{itemize}
\end{itemize}
\begin{center}\includegraphics[width=0.82\textwidth,height=0.68\textheight,keepaspectratio]{\reviewnotesimagepath d5-figure-muscle-physiology-low-atpase-activity-myosin.png}\par\captionof{figure}{The illustration relates to Low ATPase activity of myosin isoform in smooth muscle \(\rightarrow\) maximum speed of smooth muscle is slower than skeletal muscle, within the topic of Muscle Physiology.}\labfig{appendix-figure-d5-muscle-physiology-low-atpase-activity-myosin}\end{center}
\subsubsection{Cardiac Muscle Physiology}
\begin{center}
\captionof{table}{Comparison of skeletal, smooth, and cardiac muscle.}\labtab{appendix-table-d5-muscle-types}
\small
\renewcommand{\arraystretch}{1.15}
\begin{tabularx}{\textwidth}{p{0.21\textwidth}XXX}
\toprule
\textbf{Property} & \textbf{Skeletal muscle} & \textbf{Smooth muscle} & \textbf{Cardiac muscle} \\
\midrule
Shape & Cylinder & Spindle & Cylinder \\

Size & 10--100~$\mu$m & 5--10~$\mu$m & 50~$\mu$m \\

Striations & + & -- & + \\

Nucleus & Multiple & Single & Single \\

Intercalated disc & -- & -- & + \\

Innervation & Somatic & ANS & ANS \\

Function & Mechanical work & Propel luminal content & Pump blood \\
\bottomrule
\end{tabularx}
\end{center}
\begin{itemize}
\item Cardiac cells have small size w/ single nucleus and intercalated disk
\item Cells not elongated, but same striations as skeletal muscle
\item Wide, short T tubules w/ diad junctions (not triad)
\item Cardiac cells anastamose (branch, connect to other cells)
\item \textbf{Connected by intercalated disks (w/ gap junctions) and strong desmosomes}
\item Central nuclei
\item Many mitochondria
\item Cell-cell coupling is both electrical and mechanical (gap junctions and desmosomes/adherent junctions) allowing for ion depolarization cascade
\item \textbf{Excitation-contraction coupling}
\begin{itemize}
\item Depolarizing current received from gap junction
\item Action potential travels along plasma membrane and T tubules
\item Calcium channels open in plasma membrane and sarcoplasmic reticulum once they receive depolarizing current
\item Ca2+ influx from T tubule causes sarcoplasmic reticulum plasma release (Calcium induced Calcium release from sarcoplasmic reticulum)
\item Ca2+binds to troponin, which subsequently binds to tropomyosin and opens myosin binding sites
\item Crossbridge cycle activates (heart contracts)
\item Ca2+ actively reuptake by the sarcoplasmic reticulum, and mitochondria and pushed out of cell (by sarcolemmal Na-Ca exchanger, sarcolemmal Ca2+ pump)
\item Troponin releases Ca2+, muscle relaxes
\item Cardiac glycosides increase output by disabling the Sodium-Potassium ATPase pump, resulting in deeper contractions (reducing Sodium efflux)
\item Contractility - strength of a contraction at any given preload. Substances that increase contractility (positive inotropic agents) include hormones (glucagons, thyroxine), catecholemines (epinephrine and norepinephrine), drugs (digitalis) and increased calcium concentration in the extracellular fluid
\end{itemize}
\item Heart conduction system
\begin{itemize}
\item The heart beat is self sustaining by self excitation of action potentials. \(\rightarrow\) autorhythmicity by SA-AV node
\item Action potential begins at the right atrium and moves downwards and left. Visible on electrocardiogram
\item Cardiac muscle has a long refractory period (200msec vs skeletal’s 5msec and motor neuron’s 2msec)
\item Cardiac tissue cannot have tetanus (rapid repeated stimulation leading to increasing levels of contraction) because of the long refractory period (to allow time for bloodflow)
\item The long refractory period is due to net balance in uptaking and excretion of K+ and Ca2+ being equal. It lasts as long as muscle twitch
\end{itemize}
\item Cardiac muscle has 3 tissues: atrial muscle, ventricular muscle and conductive muscle fibers
\item Syncytial nature allows mass depolarization simultaneously
\item Frank-Starling's law of the heart (length-tension relationship of cardiac muscle): increased tension (stretch) from increased ventricular volume results in stronger contraction \(\rightarrow\) allows for the pumping of varied volumes (steeper increase of force for less length than skeletal muscle, narrower optimal range)
\item Law of the heart occurs because the more stretched cardiac muscle leads to myosin fibers closer to actin = more force of contraction + stretch activated Ca2+ channels for ECF Ca2+ entry
\item The heart can be stimulated to either contract or relax at a faster rate:
\begin{itemize}
\item Inotropy (increasing rate, force of a contraction) - by L type Ca2+ channels increasing Ca2+ entry, and by Ryanolyn receptors increasing Ca2+ release intracellularly
\item Lusitropy (increasing rate of relaxation) - by Phospholamban to increase velocity of Ca2+ transport by sarcoplasmic reticulum vesicle and by Tnl, decreasing affinity for TroponinC for Ca2+
\end{itemize}
\item Heart functions under aerobic mechanism, is unable to store glycogen \(\rightarrow\) prefer fatty acid as fuel over glucose \(\rightarrow\) ketone bodies used under stress, high energy demand
\end{itemize}

\subsection{Basic Biomechanics of Joint and Gait}
\begin{itemize}
\item Biomechanics \(\rightarrow\) mechanical laws related to movement/structure of living organisms
\item Lever – rigid body that pivots at a point
\item Torque - ability of a force to produce rotational motion around an axis (fixed point)
\item Three types of levers:
\begin{itemize}
\item First class levers - pivot between load and effort. Example = atlanto occipital joint (load = weight of head, effort exerted by neck muscles)
\item Second class levers - pivot at one end followed by load in the middle and effort at end. Example = metatarsophalangeal joint when tiptoeing (body weight vs calf muscles, achilles tendon)
\item Third class levers - pivot at one end followed by effort in the middle and load at end. Example = weight training arm (biceps vs weight at end)
\end{itemize}
\item Center of gravity \& balance
\end{itemize}

\begin{itemize}
\item Center of gravity - where an object’s mass is concentrated
\item Base of support - the part of the body in contact w/ supporting surface
\item Line of gravity - imaginary line passing vertically from CG to the earth
\item Balance requires line of gravity inside the base of support
\item Arthrology = study of joints
\item Synovial joints (free movement) - 2+ bones and articular cartilage (hyaline cartilage) within a joint capsule w/ surrounding (accessory) ligaments. Joint cavity in capsule contains synovial fluid secreted by synovial membrane
\begin{itemize}
\item Gliding joints - flat surfaces sliding past each other. Examples: end of clavicle, between carpals and tarsals, between vertebrae
\item Hinge joints - angular movement in a single direction. Eg. occipital bones and atlas, elbow, knee, ankle
\item Pivot joints - only rotation. Eg. atlas and axis, proximal radius and ulna
\item Ellipsoidal joints - angular motion in 2 planes. Eg. radius w/ proximal carpals, phalanges w/ metacarpals
\item Saddle joints - angular motion w/o rotation. Eg. thumb
\item Ball and socket joints - round head within depression. Eg. shoulder, hips
\end{itemize}
\item Cartilaginous joints (some movement) - connection formed by cartilage. Examples: pubic symphysis, sterno-costal joint, intervertebral disks
\item Fibrous joints (no movement) - formed by dense CT and act as fixed sutures. Examples include the radioulnar joint, tibio-fibular joint
\item Factors for joint stability: bony structure, supporting connective tissue, joint musculature
\item Shoulder joint
\begin{itemize}
\item Ball and socket joint (contd.) - formed by the head of the humerus and glenoid cavity in the scapula. The coraco-accromial ligament prevents upwards displacement, and is a part of the fibrous capsule.
\item Fibrous capsule - prevents dislocation. Consists of a thin fibrous membrane
\begin{itemize}
\item Ligaments: glenohumeral ligament (superior/middle/inferior)
\item Accessory ligaments: coracoclavicular (conoid/trapezoid), coracoacromial, coracohumeral, acromioclavicular
\end{itemize}
\item Clinical correlations:
\begin{itemize}
\item Displacement/dislocation/subluxation
\item Stiff/frozen shoulder
\item Bursitis of subaccromial bursa
\item Synovitis (biceps)/tendinitis (supraspinatus)/tendosynovitis
\end{itemize}
\end{itemize}
\item Elbow joint
\begin{itemize}
\item Articulation joint (hinge) - trochlea + trochleal notch of ulna \& capitulum + head of radius
\item Joint capsule thin at anterior/posterior side, thick at medial and lateral side
\item Ligaments: radial collateral ligament, ulnar collateral ligament
\item Only moves = flexion and extension
\end{itemize}
\item Wrist joint
\begin{itemize}
\item Colles fracture - fracture of radius presenting w/ dinner fork sign from a fall w/ outstretched hand
\item Fracture symptoms: swelling and point tenderness at the anatomical snuffbox. Is second most common fracture
\item Ligaments: collateral, radiocarpal
\item Movements \(\rightarrow\) flexion, extension, ulnar/radial deviation
\end{itemize}
\item Hip joint
\begin{itemize}
\item Synovial joint (ball and socket) - articular surfaces are the head of the femur and acetabulum
\item Ligaments at the femur head prevent hyperextension
\item Capsular ligaments: pubofemoral, iliofemoral, ischiofemoral
\end{itemize}
\item Knee joint
\begin{itemize}
\item Deformities: genu/coxa varus (bow legs from femur angle decrease), genu/coxa valgum (knock knees from femur angle increase), genu recurvatum (hyperextension of knees)
\end{itemize}
\item Gait \& movement
\end{itemize}
\begin{center}\includegraphics[width=0.82\textwidth,height=0.68\textheight,keepaspectratio]{\reviewnotesimagepath d5-figure-joint-mechanics-regional-gait-movement.png}\par\captionof{figure}{The illustration relates to Gait \& movement, within the topic of Joint Mechanics and Regional Anatomy.}\labfig{appendix-figure-d5-joint-mechanics-regional-gait-movement}\end{center}
\begin{itemize}
\item Swing phase = from push off to heel strike ground
\item Foot drop/steppage gait - from injury to deep fibular nerve \(\rightarrow\) anterior leg muscle inaction \(\rightarrow\) lack of dorsiflexion (foot hangs down and toes scrape ground, need high leg lift)
\end{itemize}

\subsection{Skeletal Muscle Relaxants}
\begin{itemize}
\item Two types of skeletal muscle relaxants:
\begin{itemize}
\item Neuromuscular blocking drugs: nondepolarizing (isoquinoline derivatives, steroid derivatives), depolarizing (succinylcholine)
\item Spasomolytic drugs: CNS action (GABAa agonist, GABAb agonist, α2-adrenergic agonist, NMDA antagonist), peripheral action (dantrolene, botulinum toxin)
\end{itemize}
\item Neuromuscular blockers:
\begin{itemize}
\item Nondepolarizing
\begin{itemize}
\item Either isoquinoline derivatives (cistatricurium, atracurium, turbocurarine – CAT) or steroid derivatives (pancuronium, rocuronium, vecuronium – PRV)
\item Tubocurarine - found in south african plants, used as poison darts; is the first drug as muscle relaxant during surgery. Has various side effects=rarely used
\item blocks/competitively inhibits NMJ communications
\item Principle: competitive inhibition of ACh at the nicotinic receptor of NMJ at low doses. Blocks Na+ channel at high doses (cannot be reversed by cholinesterase inhibitors)
\item Reversed by cholinesterase inhibitors (neostigmine, pyridostigmine, edrophonium, sugammadex - note: SG only works for steroid derivatives, since it uses encapsulation)
\item NOTE: not all muscles equally sensitive to nondepolarizing blockers (weaker/smaller is paralyzed first)
\end{itemize}
\item Pharmacokinetics of nondepolarizing blockers
\begin{itemize}
\item highly polar, inactive orally, poor absorption in GI
\item Rapidly excreted
\item Does not cross blood brain barrier/placenta
\item Half life depends on the metabolization site. Kidney=longer half life, liver=shorter half life
\item Tubocurarine - long half life. Duration approx 1-2hr. 60\% kidney, 40\% liver
\item Steroidal muscle relaxants - metabolize into 3-hydroxyl metabolite (which still is 40-80\% potent as parent drug, but has longer half life – can cause prolonged paralysis)
\item Isoquinoline derivatives - intermediate half life (atracurium = 20-35 mins, cistracurium = 25 mins). Inactivated by esterase enzyme in plasma (not affected by liver disease)
\item Atracurium metabolizes to ludanosine. Ludanosine can cross the blood-brain barrier, cause seizures (metabolized by liver, long half life). Also causes mast cell histamine release, reduces blood pressure
\item Cistracurium converts to less ludanosine = less histamine, more commonly used and safer (including for liver disease). Also high potency, but higher cost
\end{itemize}
\item Steroid muscle relaxants
\begin{itemize}
\item Pancuronium, vercuronium = no histamine release. Excreted by hepatobiliary. Prolonged duration in pts. w/ impaired liver function
\item Either long or intermediate action: pancuronium 60-100 mins, vecuronum (3-5 mins onset), rocuronium (1-2 mins onset, commonly used). Vec, roc have 20-35 mins
\item Pancuronium is excreted in urine (10-20\% is metabolized in the liver). Used as anesthetic in surgery, not preferred for long term use
\item 3-hydroxyl metabolite of pancuroneum, vercuroneium has neuromuscular blocking activity \(\rightarrow\) paralysis
\item Vercuronium, rocuronium is contraindicated in liver disease
\end{itemize}
\item Side effects of steroid muscle relaxants
\begin{itemize}
\item Atracurium - histamine release \(\rightarrow\) block autonomic ganglia, muscarinic receptors \(\rightarrow\) bronchospasm, tachycardia, hypotension
\item Pancuronium - muscarinic receptor block \(\rightarrow\) cardiovascular effect \(\rightarrow\) tachycardia
\item New drugs (cistracurium, rocuronium, vecuronium) - less histamine, fewer cardiovascular effects
\end{itemize}
\item Nondepolarizing blockers may enhance NMJ blocking when used w/ other drugs (inhaled anesthetics, local anesthetics, calcium channel blockers, aminoglycoside antibiotics). Effect also increased in myasthenia gravis or older patients \(\rightarrow\) indicated to reduce dose
\item Depolarizing (succinylcholine)
\begin{itemize}
\item Consisted of two covalently linked ACh molecules
\item Has 2 phases of action: phase 1, 2
\item Phase I block - opens Na+ channels, causes depolarization and transient contraction (fasciculation) and subsequent flaccid paralysis \(\rightarrow\) not reversed by cholinesterase inhibitors
\item Phase II block - desensitization of muscle to signalling (resistance to depolarization) because of receptor desensitization
\end{itemize}
\item Pharmacokinetics of succynylcholine
\begin{itemize}
\item Administered as continuous IV
\item Rapid action (< 90 sec) and short duration (<10 mins)
\item Metabolized in 3-5 mins by butyrylcholinesterase (or pseudocholinesterase in liver/plasma)
\item Levels of plasma cholinesterase determine duration of action of succinylcholine
\item Plasma cholinesterase testing results in dibuchane number. Usually used for pts. w/ a history of plasma cholinesterase deficiency
\item Liver issues \(\rightarrow\) prolonged duration of action
\end{itemize}
\item Adverse effects of depolarizing blockers
\begin{itemize}
\item Increases K+ release \(\rightarrow\) hyperkalemia
\item contraindicated in burn pts./massive tissue damage/head injuries/nerve damage (due to ACh receptor upregulation)
\item Scc stimulates muscarinic, ganglion \(\rightarrow\) bradycardia (countered w/ atropine)
\item Myalgia (muscle soreness) \(\rightarrow\) counteracted by small dose of nondepolarizing muscle relaxant
\item Increased intraocular, gastric pressure \(\rightarrow\) contraindicated in penetrating eye injuries, may lead to gastric regurgitation and aspiration
\item When used w/ halothane (volatile anesthetic) may induce malignant hyperthermia. Caused by increased intracellular Ca2+ and metabolism, severe uncontrolled contractions. Counteracted by dantrolene
\end{itemize}
\item Clinical uses of neuromuscular blocking drugs: surgical anesthetic adjuvants, surgical relaxation, endotracheal intubation, ventilation control, electroconvulsive shock therapy
\end{itemize}
\end{itemize}
\begin{center}\includegraphics[width=0.82\textwidth,height=0.68\textheight,keepaspectratio]{\reviewnotesimagepath d5-figure-skeletal-muscle-relaxants-clinical-uses-neuromuscular-blocking.png}\par\captionof{figure}{The illustration relates to Clinical uses of neuromuscular blocking drugs: surgical anesthetic adjuvants, surgical relaxation, endotracheal intubation, ventilation control, electroconvulsi, within the topic of Skeletal Muscle Relaxants.}\labfig{appendix-figure-d5-skeletal-muscle-relaxants-clinical-uses-neuromuscular-blocking}\end{center}
\begin{itemize}
\item NMJ function
\begin{itemize}
\item Acetylcholine receptors \(\rightarrow\) CNS and PNS. Nicotinic in all, mascarinic in CNS, cardiac, smooth, brain and ANS
\item Nicotinic receptors are divided into two types: N1/Nm, N2/Nn. N1 in NMJ and N2 in CNS, autonomic ganglia, adrenal medulla, autonomic ganglia of NMJ
\item Muscarinic localization \(\rightarrow\) M1 in gastric parietal cell, M2 in cardiac cell, M3 in exocrine gland/smooth muscle
\item Normal process - action potential received \(\rightarrow\) ACh released \(\rightarrow\) NMJ diffusion\(\rightarrow\) motor end plate receptors receive ACh signalling \(\rightarrow\) depolarization, contraction \(\rightarrow\) ACh enzyme breaks down ligand
\item Muscle spasm - increased muscle tone w/ muscle weakness. Caused by nerve circuit disruption (of muscle control, stretch reflexes). Potential causes: spinal injury, cerebral palsy, multiple sclerosis, stroke
\end{itemize}
\item Spasmolytics
\begin{itemize}
\item Divided into either CNS or PNS action
\item Reduce excessive skeletal muscle tone w/o strength reduction
\item Reduces pain, improves mobility
\item CNS action: GABAa agonist (diazepam), GABAb agonist (baclofen), alpha2-adrenergic agonist (tizanide), NDMA antagonist (riluzole)
\item Peripheral action: Ca++ channel blocker (dantrolene), ACh release block (botulinum toxin)
\end{itemize}
\item Spasmolytic drugs
\begin{itemize}
\item Benzodiazepines
\begin{itemize}
\item interact w/ GABAa receptors in the CNS. Increases Cl-, decreases excitability \(\rightarrow\) ADR: sedation effect
\item Used to treat muscle spasms leading to trauma, antidepressant, sedative
\end{itemize}
\item Baclofen
\begin{itemize}
\item GABAb agonist \(\rightarrow\) membrane polarization via K+ conductance increase (efflux) \(\rightarrow\) reduction in Ca2+ entry \& reduced release of substance P (pain transmitter) in spinal cord
\item Clinical uses: ALS skeletal muscle spasticity, lower back pain, stiff person syndrome, trigeminal neuroglia, substance craving (alcohol, nicotine, cocaine)
\item ADRs: reduced sedation, increased seizure
\end{itemize}
\item Tizanidine
\begin{itemize}
\item Alpha-2 adrenergic agonist
\item Reduces polysynaptic reflex activity and increases presynaptic inhibition of motor neurons (stops signals from reaching muscles) \(\rightarrow\) less muscle weakness
\item Used to reduce spasticity \(\rightarrow\) stroke, multiple sclerosis
\item ADRs: drowsiness, dizziness, dry mouth, hypotension
\end{itemize}
\item Riluzole
\begin{itemize}
\item Glutamate release inhibitor
\item Inhibits voltage dependent Na+ channel in presynaptic neurons
\item Used in ALS
\end{itemize}
\item Cyclobenzodiazepines
\begin{itemize}
\item used in acute spasms resulting from muscle injury (acute local muscle spasms)
\item Unknown mechanism of action on the brain stem. Interferes w/ polysynaptic reflexes for maintaining muscle tone
\item Includes: methocarbamol, orphenadrine, carisoprodol, chlorphenesin, chlorzoxazone, metaxalone
\item Norgesic is made from orphenadrine + paracetamol
\item ADRs: sedation, antimuscarinic actions
\end{itemize}
\item Dantrolene
\begin{itemize}
\item Reduces release of calcium from the sarcoplasmic reticulum of muscle cells (by ryanodine receptor channels)
\item Used for muscle spasticity \(\rightarrow\) stroke, spinal cord injury. Also is first line for malignant hyperthermia
\item Off label treatment for neuroleptic malignant syndrome caused by antipsychotics (muscle rigidity, hyperthermia)
\end{itemize}
\item Sedation of spasmolytic drugs: diazepam > baclofen > dantrolene
\item Botulinum toxin
\begin{itemize}
\item Toxin which destroys synaptobrevin (acetylcholine exocytosis protein)
\item Used for cosmetic injections (relieve wrinkles temporarily), hyperhidrosis (excessive sweating), dystonia (uncontrolled muscle spasms), strabismus (eye misalignment), skeletal muscle spasms (hemifacial, lower esophageal, cerebral palsy)
\item Protective for chronic migraine
\item ADRs: pain at injection site, infection, inflammation, swelling, redness, muscle weakness
\end{itemize}
\end{itemize}
\end{itemize}
\begin{center}\includegraphics[width=0.82\textwidth,height=0.68\textheight,keepaspectratio]{\reviewnotesimagepath d5-figure-skeletal-muscle-relaxants-adrs-pain-at-injection.png}\par\captionof{figure}{The illustration relates to ADRs: pain at injection site, infection, inflammation, swelling, redness, muscle weakness, within the topic of Skeletal Muscle Relaxants.}\labfig{appendix-figure-d5-skeletal-muscle-relaxants-adrs-pain-at-injection}\end{center}
\subsubsection{PUQ}
\begin{enumerate}
\item Which of the following is a characteristic of nondepolarizing blockers? A. Histamine blocking action B. Stimulation of ACh release C. Block of post-tetanic potentiation D. Competing w/ ACh at nicotinic receptor E. Desensitize the end plate by occupying receptor
\item Which of the following drugs is most often associated w/ hypotension caused by histamine release? A. Cistracurium B. Vecuronium C. Rocuronium D. Neostigmine E. Tubocurarine
\item Which of the following statements is INCORRECT? A. Rocuronium has fastest onset of action (for nondepolarizing blockers) B. The breakdown product of atracurium is laudanosine C. Nondepolarizing blockers do not cross cell membrane D. Drugs that are excreted by the kidney have longer half-lives E. Sugammadex rapidly reverses nondepolarizing isoquinoline derivatives
\item Which of the following is associated w/ Succinylcholine? A. Induced histamine release B. Blockade of autonomic ganglia C. Promoted intragastric and intraocular pressure D. Metabolism at the NMJ by acetylcholinesterase E. Antagonist by neostigmine during the early phase of blockade
\item Which of the following is correct regarding neuromuscular blockers (NMBs)? A. Malignant hyperthermia can be relieved by administration of atropine B. Effects of depolarizing NMBs can be reversed using cholinesterase inhibitors C. Aminoglycoside antibiotics decrease the effect of nondepolarizing NMBs D. Prior administration of nondepolarizing NMB in the operating room is used to prevent postoperative pain caused by succinylcholine E. Administration of isoquinoline derivatives causes less histamine release from mast cells
\item A 54-year-old male w/ amyotrophic lateral sclerosis (ALS) complains of generalized muscle spasm. His physician prescribes diazepam to reduce spasticity. Which of the following best explains the mechanism of the spasmolytic effect of diazepam? A. Blockade of calcium release from the sarcoplasmic reticulum B. Blockade of acetylcholine release from motor nerves C. Depolarization blockade of Nm receptors D. Facilitation of GABAA actions in the spinal cord E. Activation of GABAB receptors in the spinal cord
\item A 37-year-old woman w/ multiple sclerosis complains of painful leg cramps. The neurologist prescribes baclofen, three tablets daily. Which adverse effect is most likely? A. Bradycardia B. Muscle pain C. Hyperkalemia D. Hypotension E. Drowsiness
\item Which drug is most effective in the management of malignant hyperthermia? A. Baclofen B. Dantrolene C. Vecuronium D. Succinylcholine E. Tizanidine
\item Which drug has caused hyperkalemia leading to cardiac arrest in patients w/ spinal cord injury? A. Baclofen B. Dantrolene C. Vecuronium D. Succinylcholine E. Tubocurarine
\item Which drug administered in the operating room is used to prevent postoperative pain caused by succinylcholine? A. Baclofen B. Dantrolene C. Vecuronium D. Tizanidine E. Tubocurarine
\item A patient undergoes a 2-hour surgical procedure under isoflurane anesthesia supplemented by a depolarizing muscle relaxant. At the end of the procedure, a low dose of atropine is administered. What is the main reason? A. Provide postoperative analgesia B. Enhance the action of pyridostigmine C. Reverse the effect of the muscle relaxant D. Prevent activation of cardiac muscarinic receptors E. Prevent spasm of gastrointestinal smooth muscle
\item A 30-year-old woman presents after a motor vehicle crash. She has multiple facial and extremity lacerations and a large open fracture of her right femur. An orthopedic surgeon plans immediate operative repair of the femur, and a plastic surgeon wants to suture the facial lacerations. Which muscle relaxant would you choose for intubation?
\item Would you choose the same agent if she had experienced a 30\% total body burn at the time of the accident?
\end{enumerate}
\subsubsection{Answers and Explanations}

\begin{enumerate}
\item D. Nondepolarizing blockers act by competing w/ ACh at nicotinic receptors on the motor end plate, preventing depolarization. They do not desensitize the end plate like depolarizing blockers.
\item E. Tubocurarine is most commonly associated w/ hypotension due to histamine release. Atracurium and mivacurium can also cause this, but tubocurarine is the classic example.
\item E. Sugammadex rapidly reverses steroidal nondepolarizing blockers (e.g., rocuronium, vecuronium) but does not rapidly reverse isoquinoline derivatives (e.g., atracurium, cisatracurium). All other statements are correct.
\item C. Succinylcholine can increase intragastric and intraocular pressure. It may also cause histamine release, but it is not metabolized at the NMJ by acetylcholinesterase and neostigmine does not antagonize it during the early phase.
\item D. Atropine does not relieve malignant hyperthermia, depolarizing NMBs are not reversed by cholinesterase inhibitors, aminoglycosides enhance rather than decrease nondepolarizing NMB effects,  nondepolarizing NMBs are used to prevent succinylcholine-induced postoperative pain.
\item D. Diazepam enhances GABAA receptor-mediated inhibition in the spinal cord, reducing motor neuron excitability and spasticity.
\item E. Baclofen can cause drowsiness, dizziness, and weakness. Hypotension may occur but is less common; bradycardia, muscle pain, and hyperkalemia are not typical.
\item B. Dantrolene is the drug of choice for malignant hyperthermia as it inhibits calcium release from the sarcoplasmic reticulum.
\item D. Succinylcholine can cause life-threatening hyperkalemia in patients w/ burns, spinal cord injury, or neuromuscular diseases.
\item C. Nondepolarizing NMBs such as vecuronium are sometimes given before succinylcholine to reduce fasciculations and postoperative myalgia.
\item D. Atropine is administered to prevent muscarinic effects (e.g., bradycardia) when acetylcholinesterase inhibitors are used for reversal.
\item Succinylcholine. Rapid onset and short duration make it ideal for emergency intubation in trauma patients.
\item No, use a nondepolarizing NMB. In burn patients (>24–48h post-injury), succinylcholine may cause severe hyperkalemia. Nondepolarizing agents (e.g., rocuronium, vecuronium) are safer.
\end{enumerate}

\subsection{Physical examination of musculoskeletal system}
\begin{itemize}
\item Approach: history taking \(\rightarrow\) physical examination \(\rightarrow\) plain radiography \(\rightarrow\) other investigation: EMG/NCV, USG, CT scan, MRI
\item Clinical history
\begin{itemize}
\item Characterize pain
\item Location of pain
\item Night pain
\item Weakness
\item Deformity
\item Instability
\item Locking / Clicking / Clunking
\item Sport / Occupation
\item Previous treatments
\item Alleviating / Exacerbating
\item Acute vs. Chronic
\item Traumatic vs. Overuse
\item History of prior injury
\end{itemize}
\item Sequence for physical examination
\begin{itemize}
\item Inspect - presence of physical deformity, shortening, swelling, wasting, lesions or scars
\item Palpate - swelling, spasm, temperature changes, local/point tenderness (bony structures, joints, soft tissues)
\item Move - active and passive motion range (compared with opposite limb), pain, crepitation, spasticity (resistance to free movement)
\item Measure - limb length, circumference of limb segment compared
\item Special tests - specific for each organ
\end{itemize}
\item Observations / inspection
\begin{itemize}
\item Front \& Back
\item Height of shoulder \& scapulae
\item Asymmetry
\item Obvious deformity
\item Ecchymosis
\item Muscle atrophy
\begin{itemize}
\item Supraspinatus
\item Infraspinatus
\item Deltoid
\item NB: Popeye deformity = long head of biceps rupture
\end{itemize}
\end{itemize}
\item Neurovascular examination
\begin{itemize}
\item Motor
\item Sensory
\item Reflexes
\item Special signs; long tract sign, cerebellar sign, etc.
\end{itemize}
\item Vascular examination
\begin{itemize}
\item Pulse
\item Capillary refill
\item Ankle-brachial index (ABI)
\item Skin condition - colour, hair, etc.
\end{itemize}
\item Shoulder examination
\begin{itemize}
\item Fracture-dislocation
\item Impingement syndromes
\item Rotator cuff tear
\begin{itemize}
\item S – Supraspinatus
\item I – Infraspinatus
\item T- Teres minor
\item S- Supscapularis
\end{itemize}
\item Biceps diseases
\item Brachial plexus injury
\end{itemize}
\item Strength testing \(\rightarrow\) compare muscle group between sides
\begin{itemize}
\item Grade strength on 0 \(\rightarrow\) 5 scale
\item 0: no contraction
\item 1: muscle flicker; no movement
\item 2: motion, but not against gravity
\item 3: motion against gravity, but not resistance
\item 4: motion against resistance
\item 5: normal strength
\end{itemize}
\item Special tests
\begin{itemize}
\item Percussion Test (fracture)
\item Cozen’s Resistive Tennis Elbow Test (tendinitis)
\item Golfer’s Elbow Test (tendinitis)
\item Tinneal’s sign (nerve sensitivity), phalen’s test \(\rightarrow\) carpal tunnel syndrome
\end{itemize}
\end{itemize}
\begin{center}\includegraphics[width=0.82\textwidth,height=0.68\textheight,keepaspectratio]{\reviewnotesimagepath d5-figure-physical-examination-musculoskeletal-tinneal-sign-nerve-sensitivity.png}\par\captionof{figure}{The illustration relates to Tinneal’s sign (nerve sensitivity), phalen’s test \(\rightarrow\) carpal tunnel syndrome, within the topic of Physical examination of musculoskeletal system.}\labfig{appendix-figure-d5-physical-examination-musculoskeletal-tinneal-sign-nerve-sensitivity}\end{center}
\begin{itemize}
\item Test for de quervain's syndrome (dorsal compartment sheath inflammation)
\item Common knee problems
\begin{itemize}
\item Fracture-dislocation
\item Ligamentous injury: ACL, PCL, MCL, LCL,PLC
\item Other tissue injuries: meniscus, cartilage
\item Osteoarthritis
\item Septic arthritis
\end{itemize}
\item Disk herniation \(\rightarrow\) straight leg raising test
\end{itemize}


\section{Musculoskeletal Disease, Imaging, and Care}

\subsection{Pathology of Musculoskeletal System}
\subsubsection{Bone Disease}
\begin{itemize}
\item Osteogenesis imperfecta:
\begin{itemize}
\item Disruption of bone formation due to a lack of type I collagen
\item Also affects other tissues rich in type I collagen (eye blue+transparent sclera, middle ear hearing loss, abnormal teeth formation)
\item Autosomal dominant
\item Most common CT disorder
\item Presents as reduced bone formation and skeletal fragility
\end{itemize}
\item Osteoporosis
\begin{itemize}
\item Severe osteopenia (reduction in bone mass), Deteriorated microarchitectural bone structure
\item Increased risk of bone fracture
\item Most common forms: senile and postmenopausal
\end{itemize}
\item Fractures
\begin{itemize}
\item Defined as discontinuity of bone
\item Phases of fracture healing: inflammatory phase (1-2 days) \(\rightarrow\) reparative phase (1-2 weeks) \(\rightarrow\) remodelling phase (weeks)
\end{itemize}
\item Osteonecrosis
\begin{itemize}
\item Death of bone and bone marrow w/o infection
\item Common causes: trauma, emboli, corticosteroid administration
\item Presents as asymptomatic or joint pain
\end{itemize}
\item Osteomyelitis
\begin{itemize}
\item Infection of bone leading to inflammation
\item Most commonly Staphylococcus spp.
\item Gain entry by bloodstream (hematogenous) or direct penetration
\item Most common sites: metaphyses of the long bones (knee, ankle, hip)
\item Symptoms: fever, localized pain
\end{itemize}
\item Bone tumors: benign, malignant
\begin{itemize}
\item Benign bone tumor: osteochondroma
\begin{itemize}
\item Most common benign tumor
\item Metaphyseal region of long bones, especially near the knee
\item Slow-growing lesions
\end{itemize}
\item Malignant bone tumors: Metastatic tumors - most common, Osteosarcoma, Chondrosarcoma, Ewing sarcoma
\item Metastatic tumors
\begin{itemize}
\item Most common type
\item Caused by metastasis from other tumors
\item Mostly found in adults
\item Example cancers: carcinomas of the prostate, breast, lung, thyroid, and kidney
\end{itemize}
\item Osteosarcoma
\begin{itemize}
\item Most common matrix producing bone tumor
\item Painful, progressively enlarging masses
\item most common site = distal femur
\item Radiographs: destructive and infiltrative lesions
\item Codman triangle - breaks through the cortex and lifts the periosteum
\item 10 - 20\% pulmonary metastasis at time of diagnosis
\end{itemize}
\item Chondrosarcoma
\begin{itemize}
\item Second most common malignant matrix-producing bone tumor
\item Found in age >60 yrs
\item Location: central portion of the skeleton and distal extremities
\item Symptoms: painful, progressively enlarging masses
\item Ring shaped and popcorn like calcifications
\end{itemize}
\item Ewing sarcoma
\begin{itemize}
\item Primitive neuroectodermal tumor found in <20 yrs
\item 5\% of bone tumors
\item Small, uniform, round cells
\item Location: long bones, especially the humerus, tibia, and femur
\item Clinical features: pain and swelling of the affected area
\item Bad prognosis
\end{itemize}
\end{itemize}
\end{itemize}
\subsubsection{Joint Diseases}
\begin{itemize}
\item Osteoarthritis
\begin{itemize}
\item Slow progressive articular cartilage degeneration affecting weight bearing cartilage in older individuals
\item Occurs because of articular cartilage degeneration from repeated stress
\item Chronic and low level inflammation is noted
\item Clinical features:
\begin{itemize}
\item Often found in age >80
\item Mostly affecting weight bearing synovial joints (knees. hips, vertebrae)
\item Presenting symptoms: joint pain, progressive decrease in range of movement
\item Treated mainly w/ anti-inflammatory drugs or total replacement
\end{itemize}
\end{itemize}
\item Gout
\begin{itemize}
\item Presenting as acute and painful arthritis
\item Joint space urate crystal deposition (common in big toe). Recurrent attacks \(\rightarrow\) tophi (deposits of urate crystals), joint damage \(\rightarrow\) often 1st metaphalangeal joint
\item Pseudogout \(\rightarrow\) calcium pyrophosphate deposition, pseudogout. Can mimic gout or osteoarthritis \(\rightarrow\) often knee
\end{itemize}
\end{itemize}
\subsubsection{Benign Soft Tissue Tumors}
\begin{itemize}
\item Lipoma \(\rightarrow\) most common soft tissue tumor. Benign tumor of adipose tissue, soft, mobile and painless subcutaneous mass
\item Hemangioma \(\rightarrow\) benign blood vessel tumor w/ reddish or bluish lesion. Common in skin and liver
\item \textbf{Schwannoma \(\rightarrow\) nerve sheath tumor. Often encapsulated and slow growing. Dense and loose areas, verocay bodies (tight rows separated by acellular material)}
\end{itemize}

\subsection{Nonsteroidal Anti-inflammatory and Analgesic Drugs}
\begin{itemize}
\item \textbf{Acetaminophen (paracetamol) - prescription for strain/sprain pain in children}
\item \textbf{Aspirin (Acetylsalicylic acid)}
\begin{itemize}
\item Analgesic, anti-inflammation \(\rightarrow\) inhibition of production of prostaglandin (PGE2)
\item antipyretic
\item Mechanism of action \(\rightarrow\) inhibit both COX isoforms (decrease prostaglandin synthesis)
\item Aspirin irreversibly inhibits COX (other NSAIDS reversible)
\end{itemize}
\item \textbf{Aspirin OD \(\rightarrow\) vomiting, bicarbonate (counteract metabolic acidosis), fluids, electrolytes, others as needed}
\item \textbf{Acetaminophen OD \(\rightarrow\) N-Acetylcysteine (NAC) reduces hepatotoxicity}
\end{itemize}

\subsection{Metabolic Bone Diseases}
\begin{itemize}
\item Bone changes
\begin{itemize}
\item Bone mass, mineral density varies w/ age. Bone mass increases fastest in teenagers. Peak bone mass at approx 30 y/o
\item Gender influences bone mass \(\rightarrow\) males = faster bone mass acquisition and slower decrease
\item Menopause = higher susceptibility to bone mass decline
\item Bone mass also declines slower w/ aerobic exercise
\item Bone fracture incidence increases w/ age
\end{itemize}
\item Bone growth and bone remodeling
\begin{itemize}
\item Endochondral ossification - bone growth from an original cartilage model. Used to increase length (found most in children). Requires a growth plate, w/ proliferation on the far front and resorption on the near end
\item Intramembranous ossification - osteoblast proliferation based bone synthesis to increase width of bone. Matrix deposition from osteoblasts
\item Bone remodeling - process of repairing wear damage caused by use of bone, eg. microcracks. Also contributes calcium to plasma calcium balance
\item Process of bone remodelling - calcium needed in circulation OR defects in bone detected \(\rightarrow\) local factors released by osteoblasts \(\rightarrow\) induction of osteoclast precursor \(\rightarrow\) osteoclast precursor fusion \(\rightarrow\) bone resorption \(\rightarrow\) formation of bone pit \(\rightarrow\) osteoblast lays down new matrix \(\rightarrow\) matrix mineralization, hardening
\item RANKL \(\rightarrow\) ligand expressed by osteoblasts, binds to and activates RANK. Activated complex promotes osteoclast formation, function, survival. Receptor stimulates production of M-CSFR
\item Osteoprogerin (OPG) \(\rightarrow\) natural competitor of RANKL, secreted by osteoblasts. Acts to inhibit RANK activation. Decreases bone remodeling by disabling osteoclasts
\item M-CSFR - colony stimulating factor for monocytes (later transformed into osteoosteoclast). Controlled by RANK and RANKL complex
\item Bone resorption by osteoclast \(\rightarrow\) requires sealing to contain the area of acid+enzyme. Uses integrin at borders
\end{itemize}
\end{itemize}
\begin{center}\includegraphics[width=0.82\textwidth,height=0.68\textheight,keepaspectratio]{\reviewnotesimagepath d5-figure-metabolic-bone-diseases-bone-resorption-osteoclast-requires.png}\par\captionof{figure}{The illustration relates to Bone resorption by osteoclast \(\rightarrow\) requires sealing to contain the area of acid+enzyme. Uses integrin at borders, within the topic of Metabolic Bone Diseases.}\labfig{appendix-figure-d5-metabolic-bone-diseases-bone-resorption-osteoclast-requires}\end{center}
\begin{itemize}
\item Bone resorption requires:
\begin{itemize}
\item Proton pump (ATP driven) to maintain acidic environment
\item Enzyme phosphatase to resorb bone
\item Pro-resorption signalling (see above)
\end{itemize}
\item Whole-body calcium balance / integrated hormonal actions for calcium homeostasis
\end{itemize}
\begin{center}
\captionof{table}{Calcium concentrations in body compartments.}\labtab{appendix-table-d5-calcium-concentrations}
\small
\renewcommand{\arraystretch}{1.15}
\begin{tabularx}{\textwidth}{p{0.37\textwidth}X}
\toprule
\textbf{Compartment} & \textbf{Concentration} \\
\midrule
Total serum calcium & 8.5--10.5~mg/dL (2.1--2.6~mmol/L) \\

Ionized calcium & 4.4--5.2~mg/dL (1.1--1.3~mmol/L) \\

Protein-bound calcium & 4.0--4.6~mg/dL (0.9--1.1~mmol/L) \\

Complexed calcium & 0.7~mg/dL (0.18~mmol/L) \\

Intracellular free calcium & 0.00018~mmol/L \\
\bottomrule
\end{tabularx}
\end{center}
\begin{center}\includegraphics[width=0.82\textwidth,height=0.68\textheight,keepaspectratio]{\reviewnotesimagepath d5-figure-metabolic-bone-diseases-whole-body-calcium-balance.png}\par\captionof{figure}{Calcium fluxes between the gut, extracellular fluid, cells, bone, and kidneys.}\labfig{appendix-figure-d5-metabolic-bone-diseases-whole-body-calcium-balance}\end{center}
\begin{itemize}
\item Regulated by: absorption at small intestine, kidney filtration and bone
\item Regulation occurs for ionized (free) calcium ONLY
\item Calcitonin stops bone resorption by osteoclasts. Only released in hypercalcemic crisis
\item Parathyroid hormone stimulates bone resorption to increase calcium to normal levels. Apical calcium channels are regulated by parathyroid hormone
\item Ionized calcium is transported into the bloodstream by the combined actions of:
\begin{itemize}
\item Calcium channels in the lumen
\item Calbindin in the cell
\item The sodium-calcium pump/calcium ion pump in the basolateral membrane neighboring blood vessels
\end{itemize}
\item Thiazide diuretics promote calcium reabsorption by inhibiting the Na-Cl cotransporter (also increasing urine water content, water follows Na) in the apical membrane and reducing the basolateral Na gradient, inducing the Na-Ca exchanger to function more effectively \(\rightarrow\) thiazide diuretics are indicated in hypertensive pts w/ osteoporosis, contraindicated w/ hypercalcemia
\item Parathyroid gland \(\rightarrow\) secretes parathyroid hormone in response to (only) reduced serum ionized calcium (Ca sensing receptor inhibits PTH release by inhibiting exocytosis of PTH)
\item Response to PTH \(\rightarrow\) osteoclasts increase in function (bone resorption), kidneys increase Ca reabsorption, phosphate excretion, conversion of vitamin D into calcitriol
\item Calcitriol functions to increase calcium absorption in the intestines by binding to vitamin D receptors, promoting transcription of calcium transporter proteins (TRPV5, TRPV6 for Ca entry/calbindin for calcium transport/Na-Ca and H-Ca transporter upregulation for export)
\item Osteoblast PTH receptors function to regulate OPG, RANK ligands. Presence of PTH induces decreases in OPG and increase in RANK (osteoclast proliferation)
\item Fibroblast growth factor (FGF23) \(\rightarrow\) produced by osteocytes in response to high phosphate (eg. CKD) and vitamin D (w/ parathyroid hormone). Decreases phosphate transporters, promotes phosphate excretion at proximal renal tubules
\item FGF23 also counteracts PTHs vitamin D synthesis activity at proximal renal tubules
\end{itemize}
\begin{center}\includegraphics[width=0.82\textwidth,height=0.68\textheight,keepaspectratio]{\reviewnotesimagepath d5-figure-metabolic-bone-diseases-fgf23-also-counteracts-pths.png}\par\captionof{figure}{The illustration relates to FGF23 also counteracts PTHs vitamin D synthesis activity at proximal renal tubules, within the topic of Metabolic Bone Diseases.}\labfig{appendix-figure-d5-metabolic-bone-diseases-fgf23-also-counteracts-pths}\end{center}
\begin{itemize}
\item Hypocalcemia \(\rightarrow\) tetany, Chvostek sign (unilateral facial twitch on nerve stimulation on zygomatic arch), Trousseau sign (carpopedal/hand spasm), QT prolongation, seizures (causes: accidental parathyroid excision)
\item Hypercalcemia \(\rightarrow\) renal stones (calcium oxalate crystals), bone pain (from pseudogout, cystic bone lesions), QT shortening, anxiety and altered mental status (causes: adenoma, carcinoma of parathyroid)
\item Hypoalbuminemia (e.g. CKD) \(\rightarrow\) total serum calcium is reduced, as there is less albumin to act as a calcium carrier. Free calcium remains the same
\item Respiratory alkalosis (hyperventilation) \(\rightarrow\) reduction in H+ ions leads to more calcium (replacing protons), leading to less ionized calcium \(\rightarrow\) secretion of PTH
\end{itemize}
\begin{center}\includegraphics[width=0.82\textwidth,height=0.68\textheight,keepaspectratio]{\reviewnotesimagepath d5-figure-metabolic-bone-diseases-respiratory-alkalosis-hyperventilation-reduction.png}\par\captionof{figure}{The illustration relates to Respiratory alkalosis (hyperventilation) \(\rightarrow\) reduction in H+ ions leads to more calcium (replacing protons), leading to less ionized calcium \(\rightarrow\) secretion of PTH, within the topic of Metabolic Bone Diseases.}\labfig{appendix-figure-d5-metabolic-bone-diseases-respiratory-alkalosis-hyperventilation-reduction}\end{center}
\begin{itemize}
\item Chronic kidney disease progression (CKD) \(\rightarrow\) FGF23 initially responds to high phosphate, followed by PTH. Vitamin D decreases throughout, decreasing calcium absorption (due to nephron damage) \(\rightarrow\) hypocalcemia \(\rightarrow\) hyperparathyroidism, bone loss
\item Hypercalcemia can be caused by: vitamin D excess (resulting in higher absorption), excessive osteoclast activity, excessive tubular reabsorption in kidneys, cancer (of bone, lymphoma, PTH related protein secreting tumors, ectopic hyperparathyroidism)
\item Osteoporosis
\begin{itemize}
\item Role of estrogen in bone remodeling
\begin{itemize}
\item Promotes osteoblast differentiation (from mesenchymal stem cells)
\item Inhibits osteoclast differentiation from hematopoietic stem cells
\item Inhibits sclerostin action (sclerostin inhibits bone formation), and RANKL action
\end{itemize}
\item Menopausal estrogen decline \(\rightarrow\) postmenopausal osteoporosis
\begin{itemize}
\item Osteoporosis is diagnosed based on bone mass density T score (normal >= -1)
\end{itemize}
\end{itemize}
\end{itemize}
\begin{center}
\captionof{table}{WHO T-score criteria for diagnosis of osteoporosis.}\labtab{appendix-table-d5-osteoporosis-t-score}
\small
\renewcommand{\arraystretch}{1.15}
\begin{tabularx}{0.78\textwidth}{p{0.43\textwidth}X}
\toprule
\textbf{Diagnosis} & \textbf{T-score} \\
\midrule
Normal bone mineral density & $\geq -1$ \\

Osteopenia & Between $-1$ and $-2.5$ \\

Osteoporosis & $\leq -2.5$ \\

Severe osteoporosis & $\leq -2.5$ with a fragility fracture \\
\bottomrule
\end{tabularx}
\end{center}
\noindent\textit{The T-score is the number of standard deviations by which a patient’s bone mass falls above or below the mean peak bone mass of a 30-year-old adult.}
\begin{itemize}
\item Aging factors for osteoporosis:
\begin{itemize}
\item Vitamin D deficiency (impaired renal synthesis) \(\rightarrow\) impaired calcium absorption
\item Reduced physical activity \(\rightarrow\) sarcopenia (muscle, bone mass loss) from loss of mechanical stimulation
\item Secondary hyperparathyroidism (from reduced Ca absorption) \(\rightarrow\) increased bone resorption
\end{itemize}
\item Rickets and osteomalacia
\begin{itemize}
\item Is a bone mineralization disorder. Leads to delayed bone growth, expanded growth plates (lack of chondrocyte apoptosis)
\item Rickets of growth plate in children
\item Marked by osteoid prevalence in adults – osteomalacia
\item Cause: hypophosphatemia w/ uncorrected hypocalcemia. Lack of vitamin D precursor causes phosphate loss (by PTH action)
\end{itemize}
\end{itemize}

\subsection{Skin, Soft Tissue and Joint Infection}
\subsubsection{Skin \& Soft Tissue Infections}
\begin{itemize}
\item LO: Common skin, soft tissue, musculoskeletal infection (mostly bacteria). Pathology (+location), findings, pathogenesis
\item Not focus treatment, medication
\item Common causative agents: bacterial, fungal, viral
\item Skin \& soft tissue = superficial or deep (determined by involvement of deep facia). Deep infection = serious
\item Superficial:
\begin{itemize}
\item Impetigo - epidermal infection
\item Erysipelas - dermal infection
\item Foliculitis - hair follicle infection
\item Cellulitis - subcutaneous tissue infection
\end{itemize}
\item Deep:
\begin{itemize}
\item Necrotising fasciitis - infection of the deep facia
\item Pyomyositis - infection of muscle tissue (pyo=pus, myositis=muscle inflammation) - commonly \emph{Staphylococcus aureus. }Often presents w/ necrotising fasciitis
\end{itemize}
\item Osteomyelitis - bone infection
\item Infective arthritis - joints
\item Most common causative agents
\begin{itemize}
\item \emph{Staphylococcus aureus }(most common) - causes pus and abscess formation
\item Streptococci - \emph{Streptococcus pyrogenase} = group A - red skin w/ pus formation - gram +ve cocci in chain colonies
\end{itemize}
\item Necrotising fasciitis - deep infection - causes gas production, ulceration of epidermis. Rapidly spreads due to location. Common causes: \emph{Streptococcus pyrogenes} (group A streptococci), \emph{Staphylococcus aureus}, mixed anaerobic+aerobic bacteria, \emph{Vibrio vulnificus} (gram -ve)
\end{itemize}
\begin{center}\includegraphics[width=0.82\textwidth,height=0.68\textheight,keepaspectratio]{\reviewnotesimagepath d5-figure-skin-soft-tissue-necrotising-fasciitis-deep-infection.png}\par\captionof{figure}{The illustration relates to Necrotising fasciitis - deep infection - causes gas production, ulceration of epidermis. Rapidly spreads due to location. Common causes: Streptococcus pyrogenes, within the topic of Skin, Soft Tissue and Joint Infection.}\labfig{appendix-figure-d5-skin-soft-tissue-necrotising-fasciitis-deep-infection}\end{center}
\begin{itemize}
\item Findings = crusted lesion, impetigo (epidermal infection). \emph{S. aureus}. Found in young children
\end{itemize}
\begin{center}\includegraphics[width=0.82\textwidth,height=0.68\textheight,keepaspectratio]{\reviewnotesimagepath d5-figure-skin-soft-tissue-findings-crusted-lesion-impetigo.png}\par\captionof{figure}{The illustration relates to Findings = crusted lesion, impetigo (epidermal infection). S. aureus. Found in young children, within the topic of Skin, Soft Tissue and Joint Infection.}\labfig{appendix-figure-d5-skin-soft-tissue-findings-crusted-lesion-impetigo}\end{center}
\begin{center}\includegraphics[width=0.82\textwidth,height=0.68\textheight,keepaspectratio]{\reviewnotesimagepath d5-figure-skin-soft-tissue-skin-soft-tissue-joint.png}\par\captionof{figure}{The illustration relates to Skin, Soft Tissue and Joint Infection, within the topic of Skin, Soft Tissue and Joint Infection.}\labfig{appendix-figure-d5-skin-soft-tissue-skin-soft-tissue-joint}\end{center}
\begin{itemize}
\item Erisipilas (dermal infection) - key = erithema (areas of redness). Lesions w/ well defined borders. Commonly on face and hands (found on occupation which handle fish, veterinarian)
\item \emph{Erysipelothrix rhusiopathiae }= Gram +ve rods
\end{itemize}
\begin{center}\includegraphics[width=0.82\textwidth,height=0.68\textheight,keepaspectratio]{\reviewnotesimagepath d5-figure-skin-soft-tissue-erysipelothrix-rhusiopathiae-gram-ve.png}\par\captionof{figure}{The illustration relates to Erysipelothrix rhusiopathiae = Gram +ve rods, within the topic of Skin, Soft Tissue and Joint Infection.}\labfig{appendix-figure-d5-skin-soft-tissue-erysipelothrix-rhusiopathiae-gram-ve}\end{center}
\begin{itemize}
\item Cellulitis - Ill defined border of lesion, staph or strep infection
\end{itemize}
\begin{center}\includegraphics[width=0.82\textwidth,height=0.68\textheight,keepaspectratio]{\reviewnotesimagepath d5-figure-skin-soft-tissue-cellulitis-ill-defined-border.png}\par\captionof{figure}{The illustration relates to Cellulitis - Ill defined border of lesion, staph or strep infection, within the topic of Skin, Soft Tissue and Joint Infection.}\labfig{appendix-figure-d5-skin-soft-tissue-cellulitis-ill-defined-border}\end{center}
\begin{center}
\captionof{table}{Anatomic variants and predispositions associated with cellulitis.}\labtab{appendix-table-d5-cellulitis-predispositions}
\scriptsize
\renewcommand{\arraystretch}{1.15}
\begin{tabularx}{\textwidth}{p{0.32\textwidth}p{0.18\textwidth}X}
\toprule
\textbf{Anatomic variant or predisposition} & \textbf{Location} & \textbf{Likely bacterial cause} \\
\midrule
Periorbital cellulitis & Periorbital & \emph{Staphylococcus aureus}, \emph{Streptococcus pneumoniae}, group A streptococci \\

Buccal cellulitis & Cheek & \emph{Haemophilus influenzae} type b \\

Cellulitis complicating body piercing & Ear, nose, umbilicus & \emph{S. aureus}, group A streptococci \\

After mastectomy with axillary node dissection & Ipsilateral upper extremity & Non-group A $\beta$-hemolytic streptococci \\

After lumpectomy with limited axillary node dissection and breast irradiation & Ipsilateral breast & Non-group A $\beta$-hemolytic streptococci \\

After saphenous vein harvest for coronary artery bypass & Ipsilateral leg & Group A or non-group A $\beta$-hemolytic streptococci \\

After radical pelvic surgery or radiation therapy & Vulva, inguinal areas, legs & Group B and group G streptococci \\

After liposuction & Thigh, abdominal wall & Group A streptococci, peptostreptococci \\

Very early postoperative wound infection & Abdomen, chest, hip & Group A streptococci \\

Injection drug use (``skin popping'') & Extremities, neck & \emph{S. aureus}, streptococci (groups A, C, F, G)* \\

Perianal cellulitis & Perineum & Group A streptococci \\
\bottomrule
\end{tabularx}
\end{center}
\noindent\footnotesize\textsuperscript{*}Other bacteria to consider, based on isolation from skin or abscesses in this setting, include \emph{Enterococcus faecalis}, viridans group streptococci, coagulase-negative staphylococci, anaerobes (including \emph{Bacteroides} and \emph{Clostridium} spp.), and Enterobacteriaceae.
\begin{itemize}
\item Foliculitis - hair follicle infection, pustular lesion. \emph{Staphylococcus aureus} = Gram +ve cocci in clusters, coagulase positive
\end{itemize}
\begin{center}\includegraphics[width=0.82\textwidth,height=0.68\textheight,keepaspectratio]{\reviewnotesimagepath d5-figure-skin-soft-tissue-foliculitis-hair-follicle-infection.png}\par\captionof{figure}{The illustration relates to Foliculitis - hair follicle infection, pustular lesion. Staphylococcus aureus = Gram +ve cocci in clusters, coagulase positive, within the topic of Skin, Soft Tissue and Joint Infection.}\labfig{appendix-figure-d5-skin-soft-tissue-foliculitis-hair-follicle-infection}\end{center}
\begin{itemize}
\item Gas gangrene (type of necrotising fasciitis) - by \emph{Clostridium perfringens }(gram +ve big rods - boxcar shaped). Intubated, blebs of gas production and bleeding. Anaerobic bacteria, often presents in patients w/ wounds/trauma
\end{itemize}
\subsubsection{Bone Infections}
\begin{itemize}
\item LO: pathology, pathological findings, pathogenesis of infection. Common causative pathogens
\item Osteomyelitis / pyogenic osteomyelitis - bone infection
\item Typically hematogenous spread (bacteria from bacteremia) or spread from adjacent infection (such as from chronic wound infection in foot/sacral pressure sore)
\item Direct inoculation - injury or trauma infection
\item Commonly infects metaphysis in children (highly vascular)
\item Pathogens:
\begin{itemize}
\item Most common = \emph{Staphylococcus aureus}
\item Some gram -ve
\item Some pts. w/ risk factors can have special pathogens (eg. salmonella in sickle cell/hematologic disease, \emph{Pseudomonas spp} for diabetic foot infection)
\end{itemize}
\item Presentation: fever, leukocytosis (increased WBC), localized pain at site of infection, erythema, swelling
\item Investigation:
\begin{itemize}
\item By Xray (but can initially be normal 2 weeks post infection)
\item Initial finding change = periosteal elevation, lytic focus w/ sclerosis (late)
\item Earliest imaging = MRI
\item Diagnosis by bone biopsy
\end{itemize}
\item Diagnosis: history, physical exam, blood culture, bone biopsy\&culture. Real practice - usually just imaging (Xray in chronic osteomyelitis, MRI for acute)
\item Pathology of osteomyelitis
\begin{itemize}
\item Bacterial infection leading to suppurative inflammation
\item Vascular insufficiency on disease spread \(\rightarrow\) ischemic necrosis areas
\item Sequestrum - necrotic area of bone at the center
\item Involukrum - new bone which surrounds sequestrum area
\end{itemize}
\item Treatment: rid of bad bone - surgical drainage, antibiotics
\item Complications of osteomyelitis:
\begin{itemize}
\item Pathological fracture
\item Intraosseus abscess (brodies)
\item Amyloidosis (abnormal protein production secondary to infection)
\item Chrionic sinus tract formation (tract between abscess to surface, allow drainage)
\item Chronic osteomyelitis - squamous cell carcinoma at site of sinus tract site
\item Rarely osteogenic sarcoma
\end{itemize}
\item Chronic osteomyelitis from tuberculosis
\begin{itemize}
\item Chronic bone pain
\item Systemic infection (fever, night sweat, weight loss)
\item Caseating granuloma formation
\item Thoracic/lumbar localization (pott disease)
\item Differs from accute/chronic osteomyelitis because others localize to limbs
\item TB is AFB + w/ pulmonary infection
\end{itemize}
\item Treatment of osteomyelitis: surgical debridement
\end{itemize}

\subsection{Imaging of Musculoskeletal System}
\subsubsection{Bones}
\begin{itemize}
\item Typical imaging tools: plain film (AP, lateral - identify disease severity), CT (identify fractures), MRI (identify extent of disease), bone scan (scintigraphy - evaluate metastatic spread of cancers)
\item Radiographic density - depends on material (metal=whitest, air=darkest)
\item Key features: joint space, surface, line, bone texture (trabecular patterns are pathology??). Typically - cortex lighter than medulla
\item Label of a xray: Pt identification, film type (R hip AP standing film), adequate exposure, (ab/normal) joint space, (no/present) cortical disruption, normal bone trabeculation
\item Compact bone - present in diaphasis of long bones (bone shaft) with bone marrow cavity which stores fat
\item Spongy (cancelous) bone - presents in epiphysis of long bones, produces white and red corpuscles. Is the inner layer
\item Types of bone:
\begin{itemize}
\item Flat bones - protect internal organs
\item Short bones - stability, some movement
\item Long bones - support weight, facilitate movement
\item Irregular bone - complex shape, protects internal organs
\item Sesamoid - protect tendons from wear
\end{itemize}
\item Scapula has special transcapular view (prevent blocking of internals by scapula plate
\item Structure of elbow - fibrous joint capsule with humerus, radius and ulna. Presents with various ligaments to contain movement and fat pads
\item Pronation of hand leads to crossing of radius and ulna near radial tuberosity
\item Wrist presents with gilula carpal arches (III, II and II from medial to distal)
\item Knee in AP view - shows femur with patella and intercondylar fossa connected to tibial medial condyle and head of fibula + articular facets of tibia
\item Knee in Lat view - shows medial and lateral condyles superimposed, patella and patellofemoral joint with intercondylar eminence at the tibial joint
\item Ankle - shows tibia, fibula and some metatarsals w/ talus and lateral malleolus and calcaneal tendon
\item Foot - shows phalanges and metatarsals, tarsals w/ various joints
\item Toe - shows phalanx, metatarsophalangeal and joints w/ sesamoid bones
\item Traumatic bone injury: fracture
\begin{itemize}
\item Displaced fracture - bone moves out of typical alignment (also can be nondisplaced)
\item Complete fracture - bone is broken entirely, a fracture passes through the whole of the bone (also can be incomplete)
\item Closed or open fracture - open fracture is broken through the skin \(\rightarrow\) medical emergency
\item Typically, pediatric bone = larger elastic deformation region \(\rightarrow\) often buckle or greenstick fracture. Adult bones are typically easier to fracture
\item Types of fracture:
\end{itemize}
\end{itemize}
Complete fractures:

\begin{center}\includegraphics[width=0.82\textwidth,height=0.68\textheight,keepaspectratio]{\reviewnotesimagepath d5-figure-imaging-musculoskeletal-complete-fractures.png}\par\captionof{figure}{The illustration relates to Complete fractures:, within the topic of Imaging of Musculoskeletal System.}\labfig{appendix-figure-d5-imaging-musculoskeletal-complete-fractures}\end{center}
Incomplete fractures (common in pediatric patients):

\begin{center}\includegraphics[width=0.82\textwidth,height=0.68\textheight,keepaspectratio]{\reviewnotesimagepath d5-figure-imaging-musculoskeletal-incomplete-fractures-common-pediatric.png}\par\captionof{figure}{The illustration relates to Incomplete fractures (common in pediatric patients):, within the topic of Imaging of Musculoskeletal System.}\labfig{appendix-figure-d5-imaging-musculoskeletal-incomplete-fractures-common-pediatric}\end{center}
\begin{itemize}
\item Causes of fracture
\begin{itemize}
\item Fatigue fracture - from repetitive force on normal bone (eg. athlete shin splint)
\item Insufficiency fracture - from normal force on diffuse pathologic bone (eg. osteoporosis, osteomalacia)
\item Pathologic fracture- fracture with normal force on a focal pathologic lesion (eg. tumor or infection)
\end{itemize}
\item Aggressiveness of tumor, differential diagnosis \(\rightarrow\) plain radiograph. Bone resolution, sensitivity > 50\% for trabecular destruction
\item MRI - determine tumor extension
\item Radiograph reveals either lytic or sclerotic lesion (either well or poorly defined lytic), location (long bone - axial body or arch of spine) and miscellaneous (solitary-multiple, periosteal reaction/cortical destruction
\end{itemize}
\begin{center}\includegraphics[width=0.82\textwidth,height=0.68\textheight,keepaspectratio]{\reviewnotesimagepath d5-figure-imaging-musculoskeletal-radiograph-reveals-either-lytic.png}\par\captionof{figure}{The illustration relates to Radiograph reveals either lytic or sclerotic lesion (either well or poorly defined lytic), location (long bone - axial body or arch of spine) and miscellaneous, within the topic of Imaging of Musculoskeletal System.}\labfig{appendix-figure-d5-imaging-musculoskeletal-radiograph-reveals-either-lytic}\end{center}
\begin{itemize}
\item Features determining aggressiveness of tumors:
\begin{itemize}
\item Tumor margin - lodwick classification for solitary, lucent (radiolucent/lytic) bone lesions
\begin{itemize}
\item 1 (geographic) - A=thin, sclerotic margin (bone formation around lesion), B=distinct, well defined border (non-sclerotic), C=indistinct border
\item 2 - moth eaten
\item 3 - permeative
\end{itemize}
\item Cortical destruction - loss of cortex integrity w/ underlying lysis / endosteal scalloping
\item Periosteal reaction - reaction of cortical bone to tumor. From solid to codmans (benign to very aggressive)
\item Extraosseous soft tissue formation, soft tissue calcification not connecting to bone
\item Location of lesion
\end{itemize}
\end{itemize}
\begin{center}\includegraphics[width=0.82\textwidth,height=0.68\textheight,keepaspectratio]{\reviewnotesimagepath d5-figure-imaging-musculoskeletal-location-lesion.png}\par\captionof{figure}{The illustration relates to Location of lesion, within the topic of Imaging of Musculoskeletal System.}\labfig{appendix-figure-d5-imaging-musculoskeletal-location-lesion}\end{center}
\begin{center}\includegraphics[width=0.82\textwidth,height=0.68\textheight,keepaspectratio]{\reviewnotesimagepath d5-figure-imaging-musculoskeletal-imaging-musculoskeletal-system.png}\par\captionof{figure}{The illustration relates to Imaging of Musculoskeletal System, within the topic of Imaging of Musculoskeletal System.}\labfig{appendix-figure-d5-imaging-musculoskeletal-imaging-musculoskeletal-system}\end{center}
\begin{center}\includegraphics[width=0.82\textwidth,height=0.68\textheight,keepaspectratio]{\reviewnotesimagepath d5-figure-imaging-musculoskeletal-imaging-musculoskeletal-system-02.png}\par\captionof{figure}{The illustration relates to Imaging of Musculoskeletal System, within the topic of Imaging of Musculoskeletal System.}\labfig{appendix-figure-d5-imaging-musculoskeletal-imaging-musculoskeletal-system-02}\end{center}
\begin{itemize}
\item Osteolytic lesions appear lucent (eg. metastasis, myeloma)
\item Osteoblastic metastasis appears higher density
\item Winking owl sign, absence of pedicle. Typically osteolytic. Differentials: congenital absence of pedicle, neurofibromatosis, radiation therapy
\end{itemize}
\subsubsection{Infective diseases}
\begin{itemize}
\item Osteomyelitis - most commonly through hematogenous spread
\begin{itemize}
\item In children - does not reach growth plate (lack of blood supply)
\item Either pyrogenic (acute, subacute, chronic stages) or nonpyrogenic (TB or fungi)
\end{itemize}
\item Pyogenic osteomyelitis
\begin{itemize}
\item bone lesion forms \(\rightarrow\) pus escapes to surface and forms subperiosteal abscess \(\rightarrow\) bone inside loses blood supply (subacute) \(\rightarrow\) devitalized area becomes sequestrum (separated bone) \(\rightarrow\) new bone formed beneath stripped up periosteum (involucrum) with sinuses for pus to escape = chronic osteomyelitis
\end{itemize}
\item Nonpyrogenic osteomyelitis
\begin{itemize}
\item hematogenous spread from lungs/GU tract. More common in children (small bones > long bones)
\item For adults - typically TB arthritis
\item Findings - little/no reactive sclerosis/periosteal new bone formation. Apparent soft tissue swelling, abscess
\item Fungal non-pyrogenic osteomyelitis
\item Clinical triad - painless subcutaneous mass, multiple sinus tracts, discharge w/ grains
\item Typically on foot (madura foot)
\item No specific plain film findings
\end{itemize}
\item Infective arthritis
\begin{itemize}
\item Most commonly through hematogenous spread
\item Either septic or non-pyrogenic (TB)
\end{itemize}
\item Septic arthritis
\begin{itemize}
\item Typically monoarticular at knee, hip, ankle
\item Common pathogens: S. aureus, H. influenzae (children <5 yo)
\item First finding in plain film = joint effusion, fat displacement
\item Signs of joint effusion in knee - soft tissue density in suprapatellar recess, separated fat pad in recess indistinct quadriceps border
\item Ankle joint effusion signs - teardrop shaped soft tissue, density in anterior recess
\item Loss of cortical white line (1st osseous abnormality on plain film + subchondral bone destruction), marginal bony erosion, joint space narrowing
\end{itemize}
\end{itemize}
\begin{center}\includegraphics[width=0.82\textwidth,height=0.68\textheight,keepaspectratio]{\reviewnotesimagepath d5-figure-imaging-musculoskeletal-loss-cortical-white-line.png}\par\captionof{figure}{The illustration relates to Loss of cortical white line (1st osseous abnormality on plain film + subchondral bone destruction), marginal bony erosion, joint space narrowing, within the topic of Imaging of Musculoskeletal System.}\labfig{appendix-figure-d5-imaging-musculoskeletal-loss-cortical-white-line}\end{center}
\begin{itemize}
\item Non-pyrogenic (TB arthritis)
\begin{itemize}
\item Phemister’s triad present
\begin{itemize}
\item Juxta-articular osteoporosis (reduction in bone density next to site of inflammation)
\item Marginal bony erosion
\item Relatively preserved joint space
\end{itemize}
\item Common in larger joints (knee, hip) and monoarticular
\item Typically from adjacent OM
\item Insidious onset, delayed diagnosis
\item DDx = RA (severe loss of joint space)
\item Fibrous ankylosis end result
\end{itemize}
\end{itemize}
\begin{center}
\captionof{table}{Comparison of tuberculous and pyogenic arthritis.}\labtab{appendix-table-d5-tuberculous-pyogenic-arthritis}
\small
\renewcommand{\arraystretch}{1.15}
\begin{tabularx}{\textwidth}{p{0.45\textwidth}XX}
\toprule
\textbf{Feature} & \textbf{Tuberculous arthritis} & \textbf{Pyogenic arthritis} \\
\midrule
Soft-tissue swelling & + & + \\

Osteoporosis & + & $\pm$ \\

Joint-space loss & Late & Early \\

Marginal erosions & + & + \\

Bony proliferation (sclerosis, periostitis) & $\pm$ & + \\

Bony ankylosis & $\pm$ & + \\

Slow progression & + & $-$ \\
\bottomrule
\end{tabularx}
\end{center}
\noindent\textit{+ common; $\pm$ infrequent; $-$ rare or absent.}
\begin{itemize}
\item Soft tissue infection
\begin{itemize}
\item Necrotizing fasciitis - shows subcutaneous emphysema on plain film +/- foreign bodies
\item Cellulitis - non specific soft tissue swelling (skin) +/- FB
\item Pyomyositis - same as cellulitis but for muscle
\item Superficial fascia/subcutaneous tissue - cellulitis - use some imaging
\item Deep facia (2 connected parts) - peripheral (investing) deep facia, intermuscular deep fascia (at muscle compartment) \(\rightarrow\) site of necrotizing fasciitis
\item Muscle layer - pyomyositis
\item Others - bursitis, tenosynovitis
\item Multilayer involvement often because of fascia extending from skin to muscle
\item Either defined or no defined portal of entry
\item Defined portal of entry:
\begin{itemize}
\item Skin/mucosal breach
\item Organisms/spores introduced to soft tissue, exotoxin release
\item Exotoxins cause local damage, erythema, swelling
\item Capillary occlusion from platelet-leukocyte aggregates, vascular endothelium damaged
\item Erythema and swelling are widespread, bullae (fluid filled lesions) and ecchymoses (vessel leakage) develop
\item Deeper tissue infection, necrosis in all tissue layers and large bulla form. Larger venules/arterioles damaged
\end{itemize}
\item No defined portal of entry:
\begin{itemize}
\item Nonpenetrating deep tissue injury occurs, repair stimulated, leukocyte infiltration and myogenic progenitor activation
\item Pts. with bacteremia \(\rightarrow\) organisms trafficked to site of injury (vimentin mediated)
\item Exotoxins released, venous occlusion and necrosis
\item Arteries occluded, causing deep tissue necrosis to spread to upper tissue. Bullae and ecchymoses (bruising) develop
\end{itemize}
\item Necrotising fasciitis - necrotizing soft tissue injury  (necrotizing adipositis, fasciitis, myonecrosis)
\begin{itemize}
\item 50\% of cases in extremities: perineum = fournier gangrene, trunk, submandibular = ludwig angina
\end{itemize}
\item Pyomyositis
\begin{itemize}
\item Usually from hematogenous spread
\item M/C pathogen = S. aureus
\item Usually involves single muscle
\item M/C site = quadriceps
\item Iliopsoas muscle: direct extension from TB spondylitis, UTI, or GI infection
\item Viral myositis does not progress to abscess formation
\end{itemize}
\end{itemize}
\item Spondylodiscitis
\begin{itemize}
\item Either bacterial or tuberculosis
\item Most commonly through hematogenous spread w/ primary site at pedicle (showing winking owl sign) - spreading through venous system
\item Sequential chronological changes (slower for TB/fungi and faster for bacteria, which has sclerotic change in 10-12 weeks)
\end{itemize}
\end{itemize}
\begin{center}\includegraphics[width=0.82\textwidth,height=0.68\textheight,keepaspectratio]{\reviewnotesimagepath d5-figure-imaging-musculoskeletal-sequential-chronological-changes-slower.png}\par\captionof{figure}{The illustration relates to Sequential chronological changes (slower for TB/fungi and faster for bacteria, which has sclerotic change in 10-12 weeks), within the topic of Imaging of Musculoskeletal System.}\labfig{appendix-figure-d5-imaging-musculoskeletal-sequential-chronological-changes-slower}\end{center}
\begin{center}
\captionof{table}{Imaging features of pyogenic and tuberculous spondylitis.}\labtab{appendix-table-d5-pyogenic-tuberculous-spondylitis}
\footnotesize
\renewcommand{\arraystretch}{1.15}
\begin{tabularx}{\textwidth}{p{0.34\textwidth}XX}
\toprule
\textbf{Variable} & \textbf{Pyogenic spondylitis} & \textbf{Tuberculous spondylitis} \\
\midrule
Para- or intraspinal abscess & Absent & Present \\

Abscess wall & Thick and irregular & Thin and smooth \\

Postcontrast paraspinal abnormal-signal margin & Ill-defined & Well-defined \\

Abscess with postcontrast rim enhancement & Disc abscess & Vertebral intraosseous abscess \\

Vertebral-body enhancement pattern & Homogeneous & Heterogeneous and focal \\

Involvement of vertebral bodies & $\leq 2$ vertebral bodies & Multiple vertebral bodies \\

Commonly involved region & Lumbar spine & Thoracic spine \\

Degree of disc preservation & Moderate to complete disc destruction & Normal to mild disc destruction \\

Bony destruction of more than half & Infrequent; mild to moderate & Frequent; more severe \\
\bottomrule
\end{tabularx}
\end{center}

\subsection{Medicine for Arthritis}
\subsubsection{Arthritis}
\begin{itemize}
\item Common condition, causes pain and inflammation in joints
\item Symptoms: joint stiffness, pain. Also redness, warmth, swelling, decreased range of motion in affected joints
\item 100+ types. Most common = gout (adaptive immune response), rheumatoid arthritis (autoimmune immune response, long term)
\end{itemize}
\paragraph{Gout}
\begin{itemize}
\item Type of arthritis caused by urate crystal deposition from high levels of urate. Contributors: red meat, seafood, alcohol, sugar, kidney failure, some drugs (eg. aspirin), family history (eg. uric acid transporter issue)
\item Most common in men
\item Characterized by sudden, burning pain and stiffness + swelling in joints (often big toe)
\item Attacks repeat unless treated
\item Typical uric acid homeostasis:
\begin{itemize}
\item Purines consumed in diet/from metabolism of nucleic acids
\item Purines broken down by liver into uric acid (final step = xanthine oxidase converts into final product) – in some animals, uricase breaks down uric acid further into Allantoin
\item Uric acid eliminated by kidneys (⅔ of all – filtered and reabsorbed) and intestine (⅓ of all)
\end{itemize}
\item Formation of urate crystals
\begin{itemize}
\item Hyperuricemia leads to deposition of monosodium urate crystals in joints, soft tissue (needle-shaped deposits)
\item WBC attaches and release proteins
\item Inflammation results
\item pH lowered by protein, more crystals form
\end{itemize}
\item Treatments of gout
\begin{itemize}
\item Decrease uric acid formation
\begin{itemize}
\item Xanthene oxidase inhibitors (allopurinol, febuxostat) - Firstline medications for prevention of recurrent gout
\item Dietary changes
\end{itemize}
\item Allopurinol
\begin{itemize}
\item Broken down by xanthine oxidase into alloxanthine
\item Acts as a purine analogue
\item Alloxanthine inhibits the functions of xanthine oxidase (conversion of hypoxanthine into xanthine, and xanthine into uric acid)
\item Blocks uric acid formation \(\rightarrow\) prevent gout attacks
\item ADRs: acute gout attack at start of therapy, stomach irritation (diarrhea, nausea, vomiting), abnormal liver tests, skin rash, kidney problems (swelling, reduced urine output),
\item Severe ADRs: hepatotoxicity, Stevens-Johnson syndrome, bone marrow suppression, inflammation of blood vessels
\item Drug interactions: ampicillin, thiazides, mercaptopurine, azathioprine (reduces metabolism, requires dose reduction of affected drugs)
\item Reduced doses for renal impaired patients
\item Must be discontinued if any allergy signs present (high risk of hypersensitivity \(\rightarrow\) Stevens-Johnson syndrome \(\rightarrow\) death)
\item Stevens-Johnson syndrome - target skin lesions, mucus membrane erosions. Caused by human leukocyte antigen variant. Prevalent in Asian countries
\end{itemize}
\item Febuxostat
\begin{itemize}
\item Potent selective xanthine oxidase inhibitor. Used to treat long term gout
\item Option for those unable to take allopurinol
\item Minimal adverse events, can be used w/ renal/hepatic disease
\item ADRs: increased hepatic serum enzyme levels, nausea, diarrhea, headache, rash
\item Serious ADRs: SJS, anaphylaxis
\end{itemize}
\end{itemize}
\end{itemize}
\begin{center}
\captionof{table}{Comparison of febuxostat and allopurinol.}\labtab{appendix-table-d5-febuxostat-allopurinol}
\small
\renewcommand{\arraystretch}{1.15}
\begin{tabularx}{\textwidth}{p{0.25\textwidth}XX}
\toprule
\textbf{Feature} & \textbf{Febuxostat} & \textbf{Allopurinol} \\
\midrule
Chemical structure & Non-purine structure & Purine analogue \\

Enzyme inhibition & Selective inhibitor of xanthine oxidase; does not inhibit other enzymes involved in purine/pyrimidine synthesis and metabolism & Inhibits xanthine oxidase; may inhibit other enzymes involved in purine/pyrimidine synthesis and metabolism \\

Elimination & Both hepatic and renal elimination; can be given safely in mild to moderate renal or hepatic impairment & Eliminated primarily via renal excretion; cannot be given safely to renally impaired patients \\

Dose adjustment & No dose adjustment required in mild to moderate renal or hepatic impairment & Dose adjustment required in renal impairment \\
\bottomrule
\end{tabularx}
\end{center}
\begin{itemize}
\item Decrease manifestation of gout attacks
\begin{itemize}
\item Use NSAIDS, colchicine, steroid injections to control
\item NSAIDs: Aspirin (but low-moderate doses reduce excretion) \(\rightarrow\) large doses of aspirin recommended, which has uricosuric effect (decreases reabsorption in proximal tubule)
\item Colchicine (second line agent for those unable to tolerate xanthine oxidase inhibitors): Suppresses gouty inflammation, used to treat only acute gouty arthritis /prophylaxis to reduce attack frequency (no uricosuric or analgesic effect)
\item Corticosteroids: Used as the second line for pts. w/ contraindication for NSAID, colchicine.
\end{itemize}
\item Colchicine:
\begin{itemize}
\item Works by accumulation in WBCs leading to decreased motility, mobilization, adhesion (and overall inflammation)
\item Inhibits inflammatory mechanisms, increases levels of anti-inflammatory mediators
\item Reduces joint inflammation caused by crystal deposits
\item ADRs: GI irritation (dose-limiting effect)
\item Chronic administration: hematologic (neutropenia, aplastic anemia), myopathy (muscle weakness), alopecia
\item overdose: renal damage, CNS depression, RS failure, paralysis
\item Contraindicated in pregnancy
\end{itemize}
\item Corticosteroids:
\begin{itemize}
\item Eg. prednisone, prednisolone, methylprednisolone, dexamethasone
\item Administered intraarticularly or orally
\item Acts as anti-inflammatory and immunosuppressant by inhibiting transcription of inflammatory genes
\item Requires dose limitation, alternate dosing, gradual decreases
\item Effects of corticosteroids:
\end{itemize}
\end{itemize}
\begin{center}\includegraphics[width=0.82\textwidth,height=0.68\textheight,keepaspectratio]{\reviewnotesimagepath d5-figure-medicine-arthritis-effects-corticosteroids.png}\par\captionof{figure}{The illustration relates to Effects of corticosteroids:, within the topic of Medicine for Arthritis.}\labfig{appendix-figure-d5-medicine-arthritis-effects-corticosteroids}\end{center}
\begin{center}\includegraphics[width=0.82\textwidth,height=0.68\textheight,keepaspectratio]{\reviewnotesimagepath d5-figure-medicine-arthritis-medicine-arthritis.png}\par\captionof{figure}{The illustration relates to Medicine for Arthritis, within the topic of Medicine for Arthritis.}\labfig{appendix-figure-d5-medicine-arthritis-medicine-arthritis}\end{center}
\begin{itemize}
\item Increase renal clearance
\begin{itemize}
\item Uricosuric drugs (probenecid, sulfinpyrazone, benzbromarone, high dose aspirin)
\item Blocks active transport sites in proximal tubule. Causes decreased reabsorption of uric acid, increasing excretion
\item Treats chronic gout w/ frequent recurrence, or for patients who cannot tolerate allopurinol/need more urate lowering
\item Uricosuric drugs (probenecid, aspirin) often result in decreased excretion of uric acid, penicillin, cephalosporins, methotrexate into urine
\item For uricosuric drugs, urine volume and urinary pH should be controlled: high urine volume, alkaline urine (pH > 6) to decrease risk of uric acid stones
\item Contraindicated in: history of urinary tract stone, impaired renal function, low dose aspirin consumption
\end{itemize}
\end{itemize}
\begin{center}
\captionof{table}{Drug options for acute and chronic gout.}\labtab{appendix-table-d5-acute-chronic-gout-drugs}
\small
\renewcommand{\arraystretch}{1.15}
\begin{tabularx}{\textwidth}{XX}
\toprule
\textbf{Acute gout} & \textbf{Chronic gout} \\
\midrule
NSAIDs: naproxen, indomethacin, diclofenac & Uricosuric drugs: probenecid, lesinurad, benzbromarone, sulfinpyrazone \\

Colchicine & Synthesis inhibitors: allopurinol, febuxostat \\

Corticosteroids: prednisolone & Drugs that increase urate metabolism: pegloticase, rasburicase \\
\bottomrule
\end{tabularx}
\end{center}
\begin{center}
\captionof{table}{Objectives of gout treatment and corresponding drugs.}\labtab{appendix-table-d5-gout-treatment-objectives}
\small
\renewcommand{\arraystretch}{1.15}
\begin{tabularx}{\textwidth}{p{0.43\textwidth}X}
\toprule
\textbf{Objective of treatment} & \textbf{Drugs} \\
\midrule
Relieve inflammation and pain & NSAIDs, colchicine, glucocorticoids \\

Prevent inflammatory responses & NSAIDs, colchicine \\

Inhibit urate formation & Allopurinol, febuxostat \\

Augment urate excretion & Probenecid \\
\bottomrule
\end{tabularx}
\end{center}
\begin{center}\includegraphics[width=0.48\textwidth,height=0.68\textheight,keepaspectratio]{\reviewnotesimagepath d5-figure-medicine-arthritis-gout-treatment-pathway.png}\par\captionof{figure}{Pathways of urate production and renal clearance, with treatment targets in gout.}\labfig{appendix-figure-d5-medicine-arthritis-gout-treatment-pathway}\end{center}
\paragraph{Rheumatoid Arthritis}
\begin{itemize}
\item Autoimmune disorder where joint linings are attacked, leading to pain, inflammation, swelling
\item Presenting w/ inflammatory cytokines: TNF alpha, IL6, IL1
\item Usually affects joints symmetrically, initially begins w/ few joints (frequently: wrists, hands, elbows, shoulders, knees, ankles)
\item Signs: swelling, stiffness, soreness, warmth, fatigue
\item Potentially reversible w/ treatments (to relieve pain, reduce swelling, prevent/slow joint damage, increase function)
\item Advanced RA causes irreversible damage to bones, cartilage, ligaments and other structures (is irreversible)
\item Treatment of RA:
\begin{itemize}
\item NSAIDS (nonselective COX inhibitors, COX2 inhibitors)
\item DMARDs (Disease Modifying Anti Rheumatic Drugs)
\item Glucocorticoids (anti-inflammatory drugs)
\end{itemize}
\item DMARDs
\begin{itemize}
\item Slow acting anti-rheumatic drugs. Reduce pain, swelling, inflammation
\item Improves function, and remission in 3-6 months
\item Two types: synthetic, biologic
\item Synthetic DMARDS: methotrexate, sulphasalazine, chloroquine/hydroxychloroquine, leflunomide, gold
\item Methotrexate - functions to inhibit dihydrofoliate reductase, reduce synthesis of thymidine \& purine nucleotides/cytokines \(\rightarrow\) reduce immune proliferation. The “gold standard” of DMARDS
\item Methotrexate ADRs: hepatotoxicity, bone marrow suppression, teratogenicity
\item Sulfasalazine - inhibits COX enzyme, LOX, PAF and cytokine production. ADRs: dyspepsia, rashes, bone marrow suppression
\item Chloroquine/hydroxychloroquine - antimalarial used to treat inflammatory arthritis. Anti-inflammatory and immunomodulatory effects. ADRs: irreversible retinal toxicity, corneal deposits
\item Gold salts: reduce phagocytotic activity of macrophage and PMN. PGE2 from synovial cells
\item Leflunomide - reduces dihydro-orotate dehydrogenase (DHODH), decreasing pyramidene synthesis \(\rightarrow\) results in immune inhibition. ADRs: diarrhea, nausea, alopecia, rash, TEN, hepatic toxicity. Contraindicated in pregnancy
\item Combination therapy - using more than 1 DMARD often increases treatment effect: Triple Therapy - Methotrexate, Sulfasalazine, Hydroxychloroquine/Double Therapy - Methotrexate \& Leflunomide, Methotrexate \& Sulfasalazine, Methotrexate \& Hydroxychloroquine
\item Biological DMARDS: TNF alpha inhibitors, receptor antagonist to IL1, receptor antagonist to IL-6, monoclonal antibody to CD-20, CD-80/86/28 inhibitors, JAK inhibitors
\item Biological DMARDs
\end{itemize}
\end{itemize}
\begin{center}\includegraphics[width=0.82\textwidth,height=0.68\textheight,keepaspectratio]{\reviewnotesimagepath d5-figure-medicine-arthritis-biological-dmards.png}\par\captionof{figure}{The illustration relates to Biological DMARDs, within the topic of Medicine for Arthritis.}\labfig{appendix-figure-d5-medicine-arthritis-biological-dmards}\end{center}
\begin{itemize}
\item JAK inhibitors (e.g. tofacitinib) - suppress pro-inflammatory cytokine signalling through JAK pathway by binding with the JAK ATP site
\end{itemize}
\begin{center}\includegraphics[width=0.82\textwidth,height=0.68\textheight,keepaspectratio]{\reviewnotesimagepath d5-figure-medicine-arthritis-jak-inhibitors-tofacitinib-suppress.png}\par\captionof{figure}{The illustration relates to JAK inhibitors (e.g. tofacitinib) - suppress pro-inflammatory cytokine signalling through JAK pathway by binding with the JAK ATP site, within the topic of Medicine for Arthritis.}\labfig{appendix-figure-d5-medicine-arthritis-jak-inhibitors-tofacitinib-suppress}\end{center}
\begin{center}\includegraphics[width=0.82\textwidth,height=0.68\textheight,keepaspectratio]{\reviewnotesimagepath d5-figure-medicine-arthritis-medicine-arthritis-02.png}\par\captionof{figure}{The illustration relates to Medicine for Arthritis, within the topic of Medicine for Arthritis.}\labfig{appendix-figure-d5-medicine-arthritis-medicine-arthritis-02}\end{center}

\subsection{Drugs Affecting Bone Mineralization}
\subsubsection{Background}
\begin{itemize}
\item Bone undergoes construction by osteoblasts, destruction by osteoclasts
\item Total remodelling of bone completes every 3-4 months in young adult
\item Bone remodeling cycle:
\begin{itemize}
\item Osteoclast driven resorption
\item Osteoclast apoptosis, preosteoblast formation
\item Osteoblast maturation and osteoid tissue construction \(\rightarrow\) mineralization
\item Resting osteocytes in the area
\end{itemize}
\item Bone requires mineral homeostasis (calcium, phosphate) for homeostasis. Ca and Ph are also important for general cell function
\item Ca, Ph abnormality \(\rightarrow\) bone deformity, osteoporosis
\item Bone also serves as a site of hematopoiesis
\item Bone mineral turnover in humans:
\end{itemize}
\begin{center}\includegraphics[width=0.82\textwidth,height=0.68\textheight,keepaspectratio]{\reviewnotesimagepath d5-figure-drugs-affecting-bone-bone-mineral-turnover-humans.png}\par\captionof{figure}{The illustration relates to Bone mineral turnover in humans:, within the topic of Drugs Affecting Bone Mineralization.}\labfig{appendix-figure-d5-drugs-affecting-bone-bone-mineral-turnover-humans}\end{center}
\begin{itemize}
\item Bone mineral homeostasis
\begin{itemize}
\item Intake from intestines \(\rightarrow\) storage in bone \(\rightarrow\) excretion in kidneys
\item Release of stored Ca regulated by parathyroid hormone (by chief cells in parathyroid), absorption by vitamin D, excretion/storage by calcitonin (parafollicular cells, C cells of thyroid)
\end{itemize}
\end{itemize}
% CONTENT AUDIT: This generated caption ends mid-parenthetical; verify it against the source image.
\begin{center}\includegraphics[width=0.82\textwidth,height=0.68\textheight,keepaspectratio]{\reviewnotesimagepath d5-figure-drugs-affecting-bone-release-stored-ca-regulated.png}\par\captionof{figure}{The illustration relates to Release of stored Ca regulated by parathyroid hormone (by chief cells in parathyroid), absorption by vitamin D, excretion/storage by calcitonin (parafollicular, within the topic of Drugs Affecting Bone Mineralization.}\labfig{appendix-figure-d5-drugs-affecting-bone-release-stored-ca-regulated}\end{center}
\begin{itemize}
\item Parathyroid hormone acts to increase calcium absorption through resorption and kidney generation of active vitamin D+phosphate excretion (counteracted by vitamin D)
\item Calcitonin causes bone resorption and increased Ca excretion
\item Ca requirements:
\begin{itemize}
\item Kids and teens 1-18: 1,300 mg
\item Adults 19-50: 1,000 mg
\item All adults 50 and older: 1,200 mg
\item Pregnant/breastfeeding women: 1,500 mg
\end{itemize}
\item Vitamin D
\begin{itemize}
\item Fat soluble steroid group
\item Enhances intestinal absorption of calcium, iron, magnesium, phosphate, zinc
\item Human variants: dietary D2 (ergocalciferol) from ergosterol and D3 (cholecalciferol) from skin by 7-dehydrochloresterol + UV
\item Calcitrol (active variant of D3) stimulates Ca2+ absorption in intestine
\end{itemize}
\end{itemize}
\begin{center}\includegraphics[width=0.82\textwidth,height=0.68\textheight,keepaspectratio]{\reviewnotesimagepath d5-figure-drugs-affecting-bone-calcitrol-active-variant-d3.png}\par\captionof{figure}{The illustration relates to Calcitrol (active variant of D3) stimulates Ca2+ absorption in intestine, within the topic of Drugs Affecting Bone Mineralization.}\labfig{appendix-figure-d5-drugs-affecting-bone-calcitrol-active-variant-d3}\end{center}
\begin{itemize}
\item Vit D requirements:
\begin{itemize}
\item Adults ages 19-70: 600 IU daily
\item Adults older than 70: 800 IU daily
\end{itemize}
\item Ca/Vit D intake deficit \(\rightarrow\) consumption of multivitamins or supplements
\item Calcium excess \(\rightarrow\) kidney stones, heart attacks, other heart issues
\item Vitamin D excess \(\rightarrow\) stomach upset, constipation, weakness
\end{itemize}
\begin{itemize}
\item Bone Diseases
\end{itemize}
\begin{center}\includegraphics[width=0.82\textwidth,height=0.68\textheight,keepaspectratio]{\reviewnotesimagepath d5-figure-drugs-affecting-bone-bone-diseases.png}\par\captionof{figure}{The illustration relates to Bone Diseases, within the topic of Drugs Affecting Bone Mineralization.}\labfig{appendix-figure-d5-drugs-affecting-bone-bone-diseases}\end{center}
\begin{center}\includegraphics[width=0.82\textwidth,height=0.68\textheight,keepaspectratio]{\reviewnotesimagepath d5-figure-drugs-affecting-bone-drugs-affecting-bone-mineralization-02.png}\par\captionof{figure}{Summary of drug classes used for abnormal serum calcium, osteoporosis, and secondary hyperparathyroidism in chronic kidney disease.}\labfig{appendix-figure-d5-drugs-affecting-bone-mineralization-summary}\end{center}
\subsubsection{Drugs Affecting Bone Homeostasis}
\begin{center}\includegraphics[width=0.82\textwidth,height=0.68\textheight,keepaspectratio]{\reviewnotesimagepath d5-figure-drugs-affecting-bone-drugs-affecting-bone-homeostasis.png}\par\captionof{figure}{The illustration relates to Drugs Affecting Bone Homeostasis, within the topic of Drugs Affecting Bone Mineralization.}\labfig{appendix-figure-d5-drugs-affecting-bone-drugs-affecting-bone-homeostasis}\end{center}
\begin{itemize}
\item Bisphosphonates
\begin{itemize}
\item Inhibit osteoclasts, induce osteoclast apoptosis
\item reduce bone turnover and increase blood calcium levels
\item Oral:
\begin{itemize}
\item Alendronate (daily, weekly)
\item Risedronate (daily, weekly, monthly)
\item Ibandronate (daily, monthly)
\end{itemize}
\item Intravenous
\begin{itemize}
\item Ibandronate
\item Zoledronic acid
\end{itemize}
\item Poorly absorbed (esp. when combined with food)
\item Excreted by kidneys
\item May persist in bone for months-years
\item ADRs: abdominal pain, nausea, diarrhea, vomiting, esophagitis, gastritis or peptic ulcer \(\rightarrow\) don’t lie flat 30 mins post consumption
\item Severe ADR: osteonecrosis of jaw
\item Aledronate, risedonate indications: prophylactic/treatment for postmenopausal bone loss, osteoporosis in men, treatment of malignancy related hypercalcemia
\end{itemize}
\item Calcitonin
\begin{itemize}
\item Includes: calcimar, fortical, miacalcin
\item Increases calcium, phosphate excretion
\item Inhibits bone resorption by osteoclast inhibition
\item Peptide is given by injection: subcutaneous or intramuscular
\item Nasal spray for synthetic salmon calcitonin
\item ADRs: nausea, GI discomfort, face flush/dermatitis, pruritus (itch) at site of injection, rhinitis (nose mucosa irritation)
\item Indication: 2nd line treatment for patients with osteoporosis who cannot tolerate other antiosteoporotic agent
\end{itemize}
\item Calcium
\begin{itemize}
\item Administered in oral or IV forms (calcium carbonate, phosphate, citrate)
\item Typically administered with antiresorptive agents
\item Adequate vitamin D is needed to enable absorption
\item ADRS: gastrointestinal - nausea, constipation, gas. Treated with increased fiber+water intake, exercise
\item Hypocalcemia
\begin{itemize}
\item Causes hypersensitivity of neurons to sodium by increasing membrane permeability \(\rightarrow\) progressive depolarization
\item Causes: tetany, muscle cramps, seizures
\item Severe \(\rightarrow\) resolved with slow infusion of 5-20 mL Ca (rapid can induce arrhythmia)
\item Not severe \(\rightarrow\) resolved with oral calcium 400-1200 mg/day
\end{itemize}
\end{itemize}
\end{itemize}
\begin{itemize}
\item Vitamin D
\begin{itemize}
\item 1,25-(OH)2-D3 (calcitriol) form enhances absorption of calcium
\item Calcitriol 0.25- 1 mcg daily is capable of raising serum calcium within 24-48 hr
\item 2 main forms \(\rightarrow\) calcitriol (active form, no activation required), calcifediol (needs kidney activation, does not impact liver)
\item ADRs: hypercalcemia (anorexia, nausea, weakness)
\item Drug interactions: using thiazide diuretics + calcitriol may cause hypercalcemia
\end{itemize}
\end{itemize}
\begin{center}\includegraphics[width=0.82\textwidth,height=0.68\textheight,keepaspectratio]{\reviewnotesimagepath d5-figure-drugs-affecting-bone-drug-interactions-using-thiazide.png}\par\captionof{figure}{The illustration relates to Drug interactions: using thiazide diuretics + calcitriol may cause hypercalcemia, within the topic of Drugs Affecting Bone Mineralization.}\labfig{appendix-figure-d5-drugs-affecting-bone-drug-interactions-using-thiazide}\end{center}
\begin{itemize}
\item Osteoporosis
\begin{itemize}
\item Bone disease resulting in decreased bone strength because of lowered bone mass density (osteoclast > osteoblast)
\item Results in increased risk of fracture (common cause of fractures in elderly)
\item Pathogenesis - imbalance between bone resorption and formation
\end{itemize}
\item Role of Estrogen
\begin{itemize}
\item Estrogen decreases osteocyte apoptosis, activates bone remodelling
\item Estrogen decreases osteoblast apoptosis, oxidative stress and NF-kB, driving bone maintenance
\item Osteoclast apoptosis is driven by estrogen. RANKL induced differentiation decreased by estrogen = reduced bone resorption
\item Estrogen may act through T cells to decrease osteoclast function
\item Reduced estrogen results in excess resorption \(\rightarrow\) bone loss, low bone mass density, osteoporosis
\item Estrogen decreases osteoporosis \(\rightarrow\) estrogen replacement therapy repairs bone
\begin{itemize}
\item Oral (pills, tablets), vaginal (creams, pessaries, rings) and transdermal (implants, patches, sprays, gels) methods of delivery
\end{itemize}
\item Long term use of estrogen (>5 years) increases breast, endometrial cancer risk (cell proliferation), deep venous thrombosis, gallbladder disease \(\rightarrow\) not prescribed for osteoporosis, but effective (eg. when used to treat other conditions)
\end{itemize}
\end{itemize}
\begin{center}\includegraphics[width=0.82\textwidth,height=0.68\textheight,keepaspectratio]{\reviewnotesimagepath d5-figure-drugs-affecting-bone-long-term-estrogen-years.png}\par\captionof{figure}{The illustration relates to Long term use of estrogen (>5 years) increases breast, endometrial cancer risk (cell proliferation), deep venous thrombosis, gallbladder disease \(\rightarrow\) not prescribe, within the topic of Drugs Affecting Bone Mineralization.}\labfig{appendix-figure-d5-drugs-affecting-bone-long-term-estrogen-years}\end{center}
\begin{center}\includegraphics[width=0.82\textwidth,height=0.68\textheight,keepaspectratio]{\reviewnotesimagepath d5-figure-drugs-affecting-bone-drugs-affecting-bone-mineralization.png}\par\captionof{figure}{The illustration relates to Drugs Affecting Bone Mineralization, within the topic of Drugs Affecting Bone Mineralization.}\labfig{appendix-figure-d5-drugs-affecting-bone-drugs-affecting-bone-mineralization}\end{center}
\begin{itemize}
\item Antiresorptive drugs
\begin{itemize}
\item Includes: bisphonates, calcitonin, estrogen, SERMS, RANKL antagonist
\item Prevent resorption of bone
\item Raloxifene:
\begin{itemize}
\item Oral selective estrogen receptor modulator (act like estrogen in specific tissue, block estrogen function in others) - SERM
\item Estrogenic actions on bone, antiestrogenic actions on uterus, breast
\item Prevents osteoporosis post menopause, reduces invasive breast cancer risk (in postmenopausal women with osteoporosis), in postmenopausal women with risk for breast cancer
\end{itemize}
\item Denosumab (RANK ligand antagonist):
\begin{itemize}
\item Human monoclonal antibody targeting RANK ligand
\item Delivered subcutaneous 60 mg every 6 months
\item Binds to RANK ligand and prevents binding to osteoclasts \(\rightarrow\) osteoclast inactivation
\item Contraindicated with hypocalcemia
\item ADRs: joint, muscle pain (arms, legs), osteonecrosis of jaw
\item Approved use: postmenopausal women w/ high risk of osteoporotic fractures/history of fractures who fail/cannot tolerate other therapy
\end{itemize}
\end{itemize}
\item Bone anabolic drugs - PTH analogue (teriparatide)
\begin{itemize}
\item 34 amino acid analogue of parathyroid hormone given daily
\item Small dose \(\rightarrow\) stimulate osteoblast activity and enhance bone formation
\item Treats osteoporosis in men, postmenopausal women failing antiresorptive therapy
\item Prolonged use limited to <2 years – no safety data, may cause osteosarcoma
\item High cost
\item ADRs: itching, swelling redness at site of injection, back spasms, depression, leg cramps
\item First drug to clearly increase bone formation
\end{itemize}
\item Current guideline for treatment of osteoporosis
\end{itemize}
\begin{center}\includegraphics[width=0.82\textwidth,height=0.68\textheight,keepaspectratio]{\reviewnotesimagepath d5-figure-drugs-affecting-bone-current-guideline-treatment-osteoporosis.png}\par\captionof{figure}{The illustration relates to Current guideline for treatment of osteoporosis, within the topic of Drugs Affecting Bone Mineralization.}\labfig{appendix-figure-d5-drugs-affecting-bone-current-guideline-treatment-osteoporosis}\end{center}
\begin{itemize}
\item Phosphate
\begin{itemize}
\item Present in plasma, ECF, cell membrane phospholipids, intracellular fluid, collagen, bone tissue
\item Functions
\begin{itemize}
\item Building, repairing bones/teeth
\item Nerve function, muscle contraction
\end{itemize}
\item Hyperphosphatemia effects: hypocalcemia (phosphate+calcium salt precipitation), soft tissue calcification, tetany
\end{itemize}
\item Treatment of secondary hyperparathyroidism in CKD
\begin{itemize}
\item CKD induces hyperparathyroidism by:
\begin{itemize}
\item Hyperphosphatemia resulting in reduced active vitamin D
\item Hypocalcemia from reduced absorption (vitamin D deficit)
\item Increased FGF-23 (from reduced phosphorus clearance) inhibits conversion of vitamin D
\end{itemize}
\item Cinacalcet
\begin{itemize}
\item Acts as a calcimetric (mimics action of calcium in tissues) by allosteric activation of calcium sensing receptors expressed in parathyroid
\item Activated parathyroid calcium sensing receptors inhibit PTH secretion
\item Bone resorption inhibited (as a result of decreased PTH)
\item Taken orally after meals (30,60 or 90 mg) \(\rightarrow\) tablet should not be spit, chewed, crushed
\item Dosing is patient dependant
\item ADRs: stomach upset, vomiting, diarrhea, dizziness, nausea, weakness
\end{itemize}
\item Oral phosphate binders
\begin{itemize}
\item Reduce absorption of phosphate in GI tract
\item Taken w/ meals to bind phosphate in ingested food
\end{itemize}
\end{itemize}
\end{itemize}
\begin{center}\includegraphics[width=0.82\textwidth,height=0.68\textheight,keepaspectratio]{\reviewnotesimagepath d5-figure-drugs-affecting-bone-taken-meals-bind-phosphate.png}\par\captionof{figure}{The illustration relates to Taken w/ meals to bind phosphate in ingested food, within the topic of Drugs Affecting Bone Mineralization.}\labfig{appendix-figure-d5-drugs-affecting-bone-taken-meals-bind-phosphate}\end{center}

\subsection{Arthritis}
\begin{itemize}
\item Joints
\begin{itemize}
\item Location where 2+ bones meet, and join for movement/flexibility
\item Consists of an articular cavity and synovial tissue lining
\end{itemize}
\end{itemize}
\begin{center}\includegraphics[width=0.82\textwidth,height=0.68\textheight,keepaspectratio]{\reviewnotesimagepath d5-figure-arthritis-consists-articular-cavity-synovial.png}\par\captionof{figure}{The illustration relates to Consists of an articular cavity and synovial tissue lining, within the topic of Arthritis.}\labfig{appendix-figure-d5-arthritis-consists-articular-cavity-synovial}\end{center}
\begin{center}\includegraphics[width=0.82\textwidth,height=0.68\textheight,keepaspectratio]{\reviewnotesimagepath d5-figure-arthritis-arthritis.png}\par\captionof{figure}{The illustration relates to Arthritis, within the topic of Arthritis.}\labfig{appendix-figure-d5-arthritis-arthritis}\end{center}
\begin{itemize}
\item Key features: synovium, synovial fluid, cartilage, tendons, ligaments, bursae
\item Synovium - membrane between joint capsule and cavity. Does not cover articular cartilage. Produces synovial fluid, makes up bursae
\item Articular cartilage - weight bearing surface of synovial joints. Absorbs transmitted forces, permits low friction+high velocity movements between bones, contributes to joint stability
\item Synovial fluid - 2.5 mL in normal knee. Provides nutrition to avascular structures (cartilage) and lubricates
\item Accompanying structures: tendons (muscle-bone) run through sheath with synovial lining, ligaments (bone-bone) stabilize the bones, bursae (closed sac w/ cells similar to synovium \(\rightarrow\) help tissue gliding)
\item Defining features of arthritis: inflammation of the joint in response to a trigger
\item Athralgia - painful joints
\end{itemize}
\subsubsection{Pathogenesis, Pathology of Arthritis}
\begin{itemize}
\item Events during arthritis
\begin{itemize}
\item Activation of resident synovial macrophages, fibroblasts
\begin{itemize}
\item Production of immune mediators (histamine, IL1, etc.)
\item Proteinase production \(\rightarrow\) protein breakdown
\item Increased vascular permeability
\end{itemize}
\item Complement activation \(\rightarrow\) neutrophil influx (from bloodstream)
\item Antigen presentation - adaptive immune cells (T, B) activate
\end{itemize}
\item Initiators of joint inflammation: pathogens (infectious arthritis), crystals (crystal arthritis), abnormal force (osteoarthritis), self-antigen (autoimmune rheumatoid arthritis)
\item Presentation of arthritis
\begin{itemize}
\item Deep, dull joint pain on passive/active movement
\item Stiffness, esp after rest
\item Inability to bear weight on joint
\item Constitutional symptoms (fever, fatigue)
\item Joint swelling (increased synovial fluid)
\item Redness of overlying skin
\item Heat over joint
\item Tenderness (pain on pressure application)
\item Reduced range of motion
\end{itemize}
\item Investigations for arthritis
\begin{itemize}
\item Synovial fluid analysis \(\rightarrow\) direct evidence. Often diagnostic
\item Blood test \(\rightarrow\) systemic inflammation indicator. Not specific. Usually WBC count, C-reactive protein, erythrocyte sedimentation rate
\end{itemize}
\end{itemize}
\begin{center}\includegraphics[width=0.82\textwidth,height=0.68\textheight,keepaspectratio]{\reviewnotesimagepath d5-figure-arthritis-blood-test-systemic-inflammation.png}\par\captionof{figure}{The illustration relates to Blood test \(\rightarrow\) systemic inflammation indicator. Not specific. Usually WBC count, C-reactive protein, erythrocyte sedimentation rate, within the topic of Arthritis.}\labfig{appendix-figure-d5-arthritis-blood-test-systemic-inflammation}\end{center}
\begin{itemize}
\item Diagnostic approach of arthritis
\begin{itemize}
\item Patient history, tests:
\begin{itemize}
\item Demographics: age, gender
\item Onset: acute (<6 wks) vs. chronic (\(\geq\)6 wks)
\item Number of joints: Mono (1) vs. oligo (2-3) vs. polyarthritis (\(\geq\)4)
\item Joint location
\item Extra-articular clinical clues
\item Synovial fluid examination
\item Other investigations
\end{itemize}
\item Concomitant disease:
\begin{itemize}
\item DM,HT,DLP: Gout, infective arthritis
\item CKD \(\rightarrow\) secondary gout, drug accumulation
\item Familial and social history: Autoimmune or metabolic diseases\(\rightarrow\) rheumatoid arthritis, SLE
\item Lifestyle factors: alcohol intake (gout)
\item Unsafe sex: risk of gonorrhea, HIV, HBV \(\rightarrow\) possible septic arthritis
\end{itemize}
\item Location and distribution of arthritis:
\begin{itemize}
\item Rheumatoid arthritis: upper extremities joints: hands, wrist, elbow
\item 1st MTP joint: Gout
\item Morning stiffness: Rheumatoid: >1 hour Osteoarthritis: short (< 30 minutes)
\end{itemize}
\item Precipitating factors
\begin{itemize}
\item Gout: Alcohol, High purine diet (seafood, organ meat)
\item Septic arthritis: Septicemia, other site of infection, penetrating wound, intra-articular injection
\item Gonococcal arthritis: Sexual behaviour, menstruation, pregnancy
\end{itemize}
\item Swelling of joints: Crystal or septic joint: usually marked swelling, erythema, skin discoloration, cellulitis
\item Palpation findings:
\begin{itemize}
\item Cystic consistency: recent arthritis e.g. gout, septic arthritis
\item Bony swelling: chronic change e.g. osteoarthritis
\item Joint motion and stability
\item Crepitus (sound on joint movement
\item Tophus
\end{itemize}
\item Timing:
\begin{itemize}
\item Acute (< 6 weeks): infectious arthritis, crystal-induced arthritis
\item Chronic (> 6 weeks): osteoarthritis, rheumatoid arthritis
\end{itemize}
\end{itemize}
\end{itemize}
\begin{center}\includegraphics[width=0.82\textwidth,height=0.68\textheight,keepaspectratio]{\reviewnotesimagepath d5-figure-arthritis-chronic-weeks-osteoarthritis-rheumatoid.png}\par\captionof{figure}{The illustration relates to Chronic (> 6 weeks): osteoarthritis, rheumatoid arthritis, within the topic of Arthritis.}\labfig{appendix-figure-d5-arthritis-chronic-weeks-osteoarthritis-rheumatoid}\end{center}
\begin{itemize}
\item Ascertain presence of arthritis
\item List possible diagnoses
\item Look for common patterns between diagnoses
\end{itemize}
\begin{center}\includegraphics[width=0.82\textwidth,height=0.68\textheight,keepaspectratio]{\reviewnotesimagepath d5-figure-arthritis-look-common-patterns-between.png}\par\captionof{figure}{The illustration relates to Look for common patterns between diagnoses, within the topic of Arthritis.}\labfig{appendix-figure-d5-arthritis-look-common-patterns-between}\end{center}
\begin{itemize}
\item Degenerative joint disease (osteoarthritis)
\item Inflammatory arthritis (rheumatoid)
\item Crystal deposit arthritis
\begin{itemize}
\item Monosodium urate crystal deposition in joints \(\rightarrow\) episodic inflammatory response \(\rightarrow\) joint destruction
\end{itemize}
\end{itemize}
\begin{center}\includegraphics[width=0.82\textwidth,height=0.68\textheight,keepaspectratio]{\reviewnotesimagepath d5-figure-arthritis-monosodium-urate-crystal-deposition.png}\par\captionof{figure}{The illustration relates to Monosodium urate crystal deposition in joints \(\rightarrow\) episodic inflammatory response \(\rightarrow\) joint destruction, within the topic of Arthritis.}\labfig{appendix-figure-d5-arthritis-monosodium-urate-crystal-deposition}\end{center}
\begin{itemize}
\item Epidemiology: most common inflammatory arthritis, 1-4\% prevalence. Male:female 9:1. Age 30-50 in male, >50 in female
\end{itemize}
\begin{center}
\captionof{table}{Typical clinical pattern of gout.}\labtab{appendix-table-d5-gout-typical-pattern}
\small
\renewcommand{\arraystretch}{1.15}
\begin{tabularx}{\textwidth}{p{0.34\textwidth}X}
\toprule
\textbf{Feature} & \textbf{Typical pattern} \\
\midrule
Demographics & Adult male \\

Onset & Acute, episodic \\

Number of joints & One (monoarthritis) \\

Joint location & Big toe, midfoot, ankle \\

Extra-articular clues & Subcutaneous nodules (tophi) \\

Synovial fluid & Inflammatory type; monosodium urate crystals \\

Other investigations & High serum uric acid level (>7~mg/dL) \\
\bottomrule
\end{tabularx}
\end{center}
\begin{itemize}
\item Gold standard for diagnosis = joint arthrocentesis +ve for: monosodium urate crystals, cloudy fluid, low viscosity, WBC 2,000-5,000 and -ve culture
\item Infectious (septic) arthritis
\begin{itemize}
\item Either bacterial (gonococcal or non-gonococcal), tuberculosis or fungal
\item Key pathology: aggressive and destructive arthritis
\end{itemize}
\item Non gonococcal bacterial arthritis
\begin{itemize}
\item Presents with purulent synovial fluid, WBC >50,000
\item Risk factors
\begin{itemize}
\item Prosthetic joints
\item Existing joint disease
\item Low socioeconomic status
\item Advanced age (>80 years)
\item Intravenous drug abuse
\item Alcoholism
\item Diabetes mellitus
\item Intraarticular injection
\item Skin infections, cutaneous ulcers
\end{itemize}
\item Common causative bacteria
\begin{itemize}
\item S. aureus 59\%
\item Streptococci (gr. A, B) 16\%
\item Gram neg. Bacilli 13\%
\item \emph{S. pneumoniae} 6\%
\item Polymicrobial 6\%
\end{itemize}
\end{itemize}
\end{itemize}
\begin{center}
\captionof{table}{Typical pattern of non-gonococcal septic arthritis.}\labtab{appendix-table-d5-septic-arthritis-pattern}
\small
\renewcommand{\arraystretch}{1.15}
\begin{tabularx}{\textwidth}{p{0.34\textwidth}X}
\toprule
\textbf{Feature} & \textbf{Typical pattern} \\
\midrule
Demographics & Older adults \\

Onset & Acute \\

Number of joints & One (monoarthritis) \\

Joint location & Knee, hip, ankle \\

Extra-articular clues & Co-existing infections, preceding fever \\

Synovial fluid & Inflammatory, septic type; positive Gram stain and culture \\

Other investigations & Positive blood culture \\
\bottomrule
\end{tabularx}
\end{center}
\begin{itemize}
\item Gonococcal arthritis
\begin{itemize}
\item A manifestation of disseminated gonococcal infection
\item Caused by \emph{Neisseria gonorrhoeae}
\item Risk factors: pregnancy, multiple partners, HIV
\item Presents with tenosynovitis, pustules, synovial fluid WBC >50,000
\end{itemize}
\end{itemize}
\begin{center}
\captionof{table}{Typical pattern of gonococcal arthritis.}\labtab{appendix-table-d5-gonococcal-arthritis-pattern}
\small
\renewcommand{\arraystretch}{1.15}
\begin{tabularx}{\textwidth}{p{0.34\textwidth}X}
\toprule
\textbf{Feature} & \textbf{Typical pattern} \\
\midrule
Demographics & Young adults \\

Onset & Acute \\

Number of joints & Varies from one to multiple \\

Joint location & Varies \\

Extra-articular clues & Pustules, tenosynovitis, concurrent genitourinary tract infection \\

Synovial fluid & Inflammatory, septic type; positive Gram stain and culture \\

Other investigations & Positive blood culture \\
\bottomrule
\end{tabularx}
\end{center}
\begin{itemize}
\item Rheumatoid arthritis
\begin{itemize}
\item Systemic autoimmune condition. Manifests as chronic symmetrical joint inflammation
\item Key pathology: chronic synovitis
\item Typical presentation:
\end{itemize}
\end{itemize}
\begin{center}
\captionof{table}{Typical clinical pattern of rheumatoid arthritis.}\labtab{appendix-table-d5-rheumatoid-arthritis-pattern}
\small
\renewcommand{\arraystretch}{1.15}
\begin{tabularx}{\textwidth}{p{0.34\textwidth}X}
\toprule
\textbf{Feature} & \textbf{Typical pattern} \\
\midrule
Demographics & Adult female \\

Onset & Chronic \\

Number of joints & Polyarthritis (>4) \\

Joint location & Wrists, hands \\

Extra-articular clues & Rheumatoid nodules \\

Synovial fluid & Inflammatory type, non-specific \\

Other investigations & Rheumatoid factor, anti-citrullinated protein antibody (anti-CCP), joint erosions on X-ray \\
\bottomrule
\end{tabularx}
\end{center}
\begin{itemize}
\item Osteoarthritis
\begin{itemize}
\item Degenerative joint disease, mediated by mechanical forces and genetic predisposition
\item Key pathology: cartilage breakdown, bony enlargement and low-grade inflammation
\item Presents with pain, stiffness, crepitus, instability and mild swelling
\end{itemize}
\end{itemize}
\begin{center}
\captionof{table}{Typical clinical pattern of osteoarthritis.}\labtab{appendix-table-d5-osteoarthritis-pattern}
\small
\renewcommand{\arraystretch}{1.15}
\begin{tabularx}{\textwidth}{p{0.34\textwidth}X}
\toprule
\textbf{Feature} & \textbf{Typical pattern} \\
\midrule
Demographics & Older adults \\

Onset & Chronic \\

Number of joints & Monoarthritis or oligoarthritis \\

Joint location & Knee, hip, fingers, spine \\

Extra-articular clues & None \\

Synovial fluid & Non-inflammatory type \\

Other investigations & X-ray: joint-space narrowing, osteophytes, subchondral bone sclerosis, and cysts \\
\bottomrule
\end{tabularx}
\end{center}

\subsection{Sprain, Strain, Fracture}
\subsubsection{Strain\&Sprain}
\begin{itemize}
\item Clinical approach: history \(\rightarrow\) physical exam \(\rightarrow\) investigation (labs, radiography) \(\rightarrow\) diagnosis, treatment
\item Sprain = stretched/torn ligament (bone-bone). Most common = ankle sprain
\item Strain = stretched/torn muscle or tendon (bone-muscle). Most common = back/hamstring
\item Indications for X-ray investigation
\begin{itemize}
\item Inability to bear weight for 4 steps
\item Bone tenderness along the distal 6 cm of the posterior edge of the tibia or tip of the medial malleolus
\item Bone tenderness along the distal 6 cm of the posterior edge of the fibula or tip of the lateral malleolus
\end{itemize}
\item Degrees of sprain
\begin{itemize}
\item 1st degree - stretched ligament, no tearing. Minimal swelling, tenderness and pain on weight bearing. No ankle instability. Recovery in 2 weeks
\item 2nd degree - partial tear in ligament. Moderate swelling and tenderness. Walking with substantial discomfort. 3-6 weeks of repair
\item 3rd degree - complete ligament tear, ankle instability. Marked swelling, tenderness, pain on weight bearing with loss of motion. 8-12 months for full healing
\end{itemize}
\item Sprain healing stages
\begin{itemize}
\item 0-72 hours - inflammatory phase, vascular disruption \(\rightarrow\) edema, hematoma, prostaglandin release
\item 3-14 days - proliferative phase, fibroblast proliferation \(\rightarrow\) collagen type III formation
\item 3 wks - 1 year - remodeling phase, collagen maturation and cross-linking \(\rightarrow\) gradual strength recovery (but not to full strength)
\end{itemize}
\item Sprain treatment
\begin{itemize}
\item RICE \(\rightarrow\) rest, immobilization (w/ walking boot for 10 days), ice pack (15-20 mins every 2 hrs), compression (to control swelling), elevation (above heart)
\item NSAIDs (nonsteroidal anti-inflammatory drugs)
\item Early functional rehabilitation (range of motion, proprioception exercise) \(\rightarrow\) allow return to normal activity
\item Return to normal activity at 3 wks
\end{itemize}
\end{itemize}
\subsubsection{Fracture\&Bone Healing}
\begin{itemize}
\item Fracture stability dictates the type of healing that will occur
\begin{itemize}
\item Mechanical stability governs the mechanical strain
\item Strain (stretch): measure of deformation in term of relative displacement
\item When the strain is below 2\%, primary bone healing will occur
\item When the strain is between 2\% and 10\%, secondary bone healing will occur
\end{itemize}
\item Primary bone healing \(\rightarrow\) intramembranous ossification by haversian remodeling. Occurs with rigid constructs, such as plate fixation
\item Secondary bone healing \(\rightarrow\) endochondral healing (callus formation). Occurs with non-rigid fixation
\item Stages of bone healing
\end{itemize}
\begin{center}\includegraphics[width=0.82\textwidth,height=0.68\textheight,keepaspectratio]{\reviewnotesimagepath d5-figure-sprain-strain-fracture-stages-bone-healing.png}\par\captionof{figure}{The illustration relates to Stages of bone healing, within the topic of Sprain, Strain, Fracture.}\labfig{appendix-figure-d5-sprain-strain-fracture-stages-bone-healing}\end{center}

\subsection{Compartment syndrome, Bone tumor}
\subsubsection{Bone tumors}
\begin{itemize}
\item Main types of bone tumor:
\begin{itemize}
\item Primary bone tumor - bone-forming, cartilage forming, miscellaneous
\item Secondary bone tumor - cancer that spreads away from original site. Usually age >50 yrs. Commonly breast, lung, thyroid, renal or prostate
\end{itemize}
\end{itemize}
\begin{center}\includegraphics[width=0.82\textwidth,height=0.68\textheight,keepaspectratio]{\reviewnotesimagepath d5-figure-compartment-syndrome-bone-secondary-bone-tumor-cancer.png}\par\captionof{figure}{The illustration relates to Secondary bone tumor - cancer that spreads away from original site. Usually age >50 yrs. Commonly breast, lung, thyroid, renal or prostate, within the topic of Compartment syndrome, Bone tumor.}\labfig{appendix-figure-d5-compartment-syndrome-bone-secondary-bone-tumor-cancer}\end{center}
\begin{itemize}
\item Osteoid osteoma (benign bone forming)
\begin{itemize}
\item Benign osteoblastic lesion, usually <1.5 cm
\item Fibrous-osseous tumor (nidus) surrounded by a zone of reactive bone sclerosis
\item M>F \textasciitilde{} 3:1, 70\% < 20 y
\item Sites: 80\% long bones (femur, tibia) near cortex
\item S \& S: rest pain (worse at night) relieved by aspirin or NSAIDs
\item X-ray indicates lytic nidus < 1.5 cm with sclerotic rim + intense periosteal reaction. CT finding indicates location and size
\item Histo: nidus of cellular, highly vascularized tissue, immature bone \& osteoid tissue with zone of sclerotic bone
\item Treatment: nonoperative w/ observation and NSAIDs. Surgical resection indicated in pts with progressive and severe symptoms
\item Percutaneous radiofrequency ablation results in 10-15\% recurrence rate
\item Open curettage with burr-down technique (identify nidus and using high-speed burr) is invasive option
\end{itemize}
\item Osteoblastoma (benign bone forming)
\begin{itemize}
\item Rare fibro-osseous tumor arising from bone
\item Histologic structure = osteoid osteoma (But larger size usually > 1.5 cm)
\item M>F 2:1, 10- 30 y
\item Sites: 42.5\% spinal column (posterior element*)
\item S \& S: Slowly progressive dull aching pain, partial response to
\item NSAIDs
\item May see neurologic symptoms with spine involvement
\item X-ray \& CT scan: Well-defined lytic region surrounded by sclerotic, reactive border of bone nidus > 1.5 cm
\item Histology: similar to osteoid osteoma but with more giant cells, w/ nidus of immature osteoid
\item Treatment: Extended curettage w/ marginal excision +/- bone grafting
\end{itemize}
\item Osteosarcoma (malignant bone forming)
\begin{itemize}
\item Malignant, aggressive osteogenic bone tumors
\item Commonly found in distal femur, proximal tibia
\item Pts typically young adults/children w/ rapidly progressive pain+swelling
\item Incidence 1 in 5 million per year
\item Bimodal distribution of occurrence
\begin{itemize}
\item First peek is the second decade of life (10-14 years of age is most common)
\end{itemize}
\item Second peak is in patients over 65 seen in those with Paget's disease and those who have received prior radiation
\item Radiographic findings: characteristic blastic and destructive lesion, with sun-burst or hair-on-end pattern of matrix mineralization periosteal reaction (Codman's triangle). MRI reveals extent of tumor
\item Tumor cells show significant atypia, and produce osteoid
\item Stroma cells show malignant characteristics with atypia, high nuclear to cytoplasmic ratio, and abnormal mitotic figures
\item Histologic features: primary, solitary, intramedullary (within bone marrow), and poorly differentiated, producing a predominantly bony matrix
\item Treatment
\begin{itemize}
\item Preoperative chemotherapy
\item Operation \(\rightarrow\) wide resection and reconstruction with endoprosthesis
\item Followed by adjuvant chemotherapy
\end{itemize}
\end{itemize}
\item Osteochondroma (benign cartilage forming)
\begin{itemize}
\item Most common benign tumor
\item Benign chondrogenic lesions may take the form of solitary osteochondroma, or multiple hereditary exostosisTypically present at ages 10-30 with painless slow-growing mass
\item Malignant transformation of a solitary osteochondroma to chondrosarcoma is 1\%
\item Malignant transformation of a multiple hereditary exostosis to chondrosarcoma is 10\%
\item Diagnosis is made with radiographs showing pedunculated or sessile lesions found on bone surface
\item Treatment \(\rightarrow\) observation in asymptomatic cases, surgical resection in progressive and severe pain
\item Histologic features - endochondral ossification on basal surface of hyaline cartilage, disordered cartilage and binucleate chondrocytes in lacunae
\end{itemize}
\item Endochondroma (benign cartilage forming)
\begin{itemize}
\item Benign lesions of hyaline cartilage
\item Typically occur in medullary cavity of the diaphysis or metaphysis
\item Most common tumor of small bones (hands, feet)
\item Asymptomatic presentation, often incidental or with pathological fracture
\item X-ray finding
\item Long bone: popcorn calcification, small bone: endosteal scalloping
\item Treatment:
\begin{itemize}
\item Long bone: observe (malignant change < 1\% )
\item Small bone: curettage + bone graft
\item Pathologic fracture: treat fracture first
\end{itemize}
\end{itemize}
\item Chondrosarcoma (malignant cartilage-forming)
\begin{itemize}
\item Malignant primary bone tumors composed of chondrocytes with variable degrees of malignancy
\item Most found in the pelvis, proximal femur, and proximal humerus
\item Patients typically present after the age of 40, with a painful growing mass
\item X-ray: Popcorn calcification similar appearance to enchondromas with additional cortical expansion and endosteal erosion
\item Diagnosis: biopsy with hypercellular stroma consisting of characteristic "blue-balls" of a cartilage lesion which permeate the bone trabeculae
\item High grade lesions: pleomorphic chondrocytes w/ frequent mitotic features, anaplastic chondrocytes with chondroid matrix
\item Treatment \(\rightarrow\) solely resection (chemotherapy resistant)
\end{itemize}
\item Giant cell tumor (misc. Bone tumors)
\begin{itemize}
\item Benign aggressive tumor
\item Typically found in the epiphysis of long bones
\item Most commonly at the distal femur and proximal tibia > 50\%
\item 5\% of benign bone tumors
\item Patients typically present between ages 30 and 50 years
\item Insidious onset of pain of the involved extremity with activity, at night, or at rest
\item Radiographic findings: eccentric lytic lesion at epiphyseal/metaphyseal area often extending into the distal epiphysis and borders subchondral bone
\item Chest radiograph or chest CT: pulmonary metastasis occurs up to 3\%
\item Treatment: Treatment is generally curettage, adjuvant treatment, and reconstruction depending on the location of the lesion \& staging
\end{itemize}
\end{itemize}
\subsubsection{Compartment Syndrome}
\begin{itemize}
\item Osseofascial compartment pressure rises \(\rightarrow\) perfusion decreases
\item May lead to irreversible muscle and neurovascular damage
\item May occur anywhere that skeletal muscle is surrounded by fascia
\item Most commonly occurs at leg, forearm \& hand
\item Commonly occurs in ant. and deep posterior compartment
\item Etiology \(\rightarrow\) Increase in volume within an enclosed space
\begin{itemize}
\item Fracture with bleeding (70\% of cases)
\item Crush injuries
\item Arterial injury
\item Post-ischemic edema
\item Burns
\item Bleeding disorders
\item Tight cast
\end{itemize}
\end{itemize}
\begin{center}\includegraphics[width=0.82\textwidth,height=0.68\textheight,keepaspectratio]{\reviewnotesimagepath d5-figure-compartment-syndrome-bone-tight-cast.png}\par\captionof{figure}{The illustration relates to Tight cast, within the topic of Compartment syndrome, Bone tumor.}\labfig{appendix-figure-d5-compartment-syndrome-bone-tight-cast}\end{center}
\begin{itemize}
\item Leg compartments
\begin{itemize}
\item Anterior compartment \(\rightarrow\) Dorsiflexion of foot and ankle (Tibialis anterior, extensor hallucis longus, extensor digitorum longus peroneus tertius)
\item Lateral compartment \(\rightarrow\) plantarflexion and eversion of foot (peroneus longus+brevis)
\item Deep posterior compartment \(\rightarrow\) Plantarflexion of foot and big toe, inversion of foot (tibialis posterior, flexor hallucis longus, flexor digitorum longus)
\item Superficial posterior compartment \(\rightarrow\) foot, ankle plantarflexion (gastrocnemius, soleus, plantaris)
\end{itemize}
\item Presentation of compartment syndrome
\begin{itemize}
\item Pain out of proportion for the severity of mechanism of injury
\item Pain on passive stretching of the muscle that suspected for compartment syndrome \(\rightarrow\) sensitive sign
\item Paresthesia (pins and needles) \(\rightarrow\) nerve ischemia
\item Pallor (poor circulation)
\item Poikilothermia (limb is colder)
\item Paralysis (late finding, usually not full recovery)
\item Pulselessness (late finding, amputation usually inevitable)
\end{itemize}
\item Management of compartment syndrome
\begin{itemize}
\item Raise limb above heart
\item If cast present, take off
\item If symptoms don’t resolve in 30 mins \(\rightarrow\) fasciotomy
\item Indications for fasciotomy
\begin{itemize}
\item Intra-compartment pressure > 30 mmHg in the presence of clinical findings
\item Major arterial injury persisting > 4 hours
\item Circumferential full-thickness burns
\end{itemize}
\end{itemize}
\end{itemize}

\subsection{Clinical correlation of head and neck, diseases}
\begin{itemize}
\item Nope sorry
\end{itemize}

\subsection{Clinical of degenerative spine and OA knee}
See the slides

\begin{center}\includegraphics[width=0.82\textwidth,height=0.68\textheight,keepaspectratio]{\reviewnotesimagepath d5-figure-degenerative-spine-oa-see-slides.png}\par\captionof{figure}{The illustration relates to See the slides, within the topic of Clinical of degenerative spine and OA knee.}\labfig{appendix-figure-d5-degenerative-spine-oa-see-slides}\end{center}
\begin{itemize}
\item Stenosis:
\begin{itemize}
\item classic clinical syndrome neurogenic claudication with pain radiating beyond the back to the buttocks, thighs or lower legs
\item pain is worsened with extension lumbar spine (standing or walking) and improves with flexion lumbar spine (sitting or stooping forward)
\end{itemize}
\end{itemize}

\end{document}